% Modernised reading — fonds Alexandre Grothendieck, shelfmark 148, pages 1-71.
%
% This is an INTERPRETATION, not a transcription. It re-reads the pages in
% today's notation and vocabulary, states the hypotheses the manuscript leaves
% implicit, and drops the critical apparatus so that the mathematics reads
% straight through. Everything the brackets and underlines carried moves into a
% footnote. It is held to being mathematically correct as it stands; where the
% manuscript is loose or mistaken the true statement appears and the footnote
% says what the page has. In any dispute of substance the transcription
% governs: whoever cites Grothendieck cites batch-01.fr.tex to batch-04.fr.tex.
%
% Scope: the folder is seventy-one pages in four batches (1-20, 21-40, 41-60,
% 61-71), all transcribed, read whole before any of this was written, and
% written as ONE file for the whole shelfmark. Pages with nothing of his
% (typed administrative versos 5-15 odd, listings 21 and 29) are not read;
% covers 16, 28, 39, 65 give section titles.
%
% The spine. No single argument: a working file of autumn 1983 (page 40 is
% dated « Oct. 1983 », page 1 « 29.11.83 ») whose runs all serve one project,
% the « tour de Teichmüller » — the system of mapping-class groupoids of all
% surfaces with the operations that relate them — and its reconstruction from
% the pieces of modular dimension <= 1 (relations from dimension 2). Stations:
%   1-3    groups of a surface with toric boundary: SRT, fractional twists R^I;
%   4-14   the two exceptional elliptic curves (j = 1728, j = 0) and their cells;
%   16-20  M_{0,4}, M_{1,1}: special base points, paths, T_{0,4} vs T_{1,1};
%   22-27  maps, pieces by modular dimension, trivalent maps and coverings;
%   28-38  M_{0,5}: the chart of 542 special structures and 1620 arrows, the
%          classification of 5-point systems with symmetry (read 37, 36, 38);
%   39-42  catalogue of operations, the « douze avatars » of the tower;
%   43-60  « Groupoïdes de Teichmüller » (his pp. 1-17): the groupoid S_*TD,
%          the Mumford-Deligne graph functor, graphs with « repères », the
%          elementary operations and their recapitulation 0°-7°;
%   61-64  (his pp. 18-21) the programme of a « petit théorème de dévissage »;
%   65-71  pants decompositions of maps; page 70 restarts the definitions
%          (a separate mouture, kept apart, with S for the marked points).
%
% Notation fixed for the whole folder. Oriented groupoids carry « + »
% explicitly (T^+_top), bioriented ones « ± »; the pages write T for the
% oriented one although the margin of p. 43 proposes the opposite. S_* prefix:
% rigidification points R; D suffix: divided by cutting circles Sigma. The
% marked set is R everywhere (page 70 calls it S, translated). In a graph, S
% is the vertex set, A-> the half-edges, A = A_0 ⊔ A_l (bound / free edges),
% gamma the genus weight, r the marking weight. The circle group is U, I =
% pi_0(dX) on pages 1-3. Z_a = Z twisted by the orientations omega(a).
% Elementary operations numbered as on page 60 (4° = gommage of a cutting
% circle, 5° = capping a hole); page 63 numbers the first 5°, translated.
%
% Substantive departures from the pages, each footnoted where it occurs:
%   - pp. 1-3: Z^I -> S_*T^!(X) injective and R^I ∩ S_*T^! = Z^I need every
%     component with 2g-2+nu > 0 (supplied, cf. p. 40 d)); SRT = R^I.S_*T(X,s)
%     is a product of subgroups, not a contracted product, outside the pure part;
%     s |-> phi_s is not injective (diagonal rotations), the torsor statement of
%     p. 2 is not asserted; on p. 3 the permutations must respect (n_i);
%   - p. 10: zeta = exp(2 i pi/6) needs the i;
%   - p. 17: s_1 = 48, forced by Euler (22 - 48 + 28 = 2) and by the edge count;
%   - p. 19: the bottom row holds with S~_3 the dicyclic group of order 12
%     (S |-> x, ST |-> a), its kernel N free of rank 2 and mapped isomorphically
%     onto pi_{0,3} = Gamma(2)/±1; the annotation S~_3 = SL(2,Z/4)/± is wrong:
%     that group has order 24 (it is S_4, checked by element orders);
%   - p. 23: the list sums to 17 (2+3+6+6), the figure written is 14 = 3x4+2;
%   - p. 25: trivalent maps correspond to coverings of P^1 ramified over 0, 1,
%     inf with indices dividing 3 and 2 over two of them; arbitrary three-point
%     coverings give bipartite maps (equivalence restricted);
%   - p. 30/32: sources and targets read off the chart do not always fit the
%     tally's factors (gamma^0_r, gamma^2_r, gamma^2'_r); the line of
%     gamma^1'_r must be 3x120 = 360, which the total 1020 requires;
%   - p. 34: « Est-ce le cas 10 ? » — no: mu_4 ∪ {e^{i pi/4}} has an
%     automorphism group of order 2 (checked numerically);
%   - p. 34: an anti-automorphism preserving K need not have a fixed point nor
%     be of order 2 (zeta_5/zbar on mu_5 has order 10, no fixed point); an odd
%     power of it is a reflection preserving K, which is what the
%     classification uses;
%   - p. 38: mu = 0 (the pyramidal point) is excluded on the page without
%     reason; it is the point with extra symmetry; lambda = mu^2;
%   - p. 40: 3g-3+nu >= 0 is not equivalent to the other two inequalities
%     (g = 1, nu = 0); tau_rho = id for rho in Z holds for the reglued object,
%     not for the arrow (a Dehn twist);
%   - p. 43: the signature chi exists only on connected surfaces (his margin);
%   - p. 49: R_a is a torsor only when r(a) >= 1;
%   - p. 56: « restriction de Z_a/r à Z_a/r' » is restriction to the subgroup
%     of order r' (isomorphic to Z/r'); « surmarquage » is extension along
%     Z/r -> Z/dr;
%   - p. 62: the unread sign in (**) is « - » (capping nu' holes);
%   - p. 66: alpha, beta, gamma in N needs the triangle inequalities besides
%     parity; the number of bigons is alpha+beta+gamma-3, not -2.
%
% Readings that are ours, flagged as such: the D_inf-set of « repères » as the
% flags of the polygons P_a (pp. 49-51, supported by p. 40 c)); T'_{0,4} as
% T_{0,4}/V (p. 19); the counts of p. 35 as those of the unprimed types of
% p. 30; the Riemann-Hurwitz check of the counts on p. 26.
%
% Announced and never established in the folder: the fundamental relations
% between the 1620 arrows (p. 32, programme 1)-3)); the « petit théorème de
% dévissage » (p. 61); the description of T_top(Gamma) functorially in Gamma
% (p. 62); the models of the fundamental groupoids Pi_1(X; R) (pp. 63-64).
%
% No third-party texts in the folder; nothing of RIGHTS.md applies.
%
% Pass: Opus 5.5 (claude-opus-5-5), 2026-10-10 — first pass, unchecked against the pages by a human.
\documentclass[11pt,a4paper]{article}
\input{../preamble/grothendieck}

\folder{148}
\pages{1}{71}
\dating{1983}
\watermark{Édition de démonstration\\interprétation personnelle de l'œuvre}
\foldertitle{Teichmüller (1983) : notes manuscrites (1983) — lecture modernisée du dossier entier}

\begin{document}

\section*{Résumé}

\begin{resume}
Une surface — une sphère, un tore, un bretzel à plusieurs anses, avec ou sans
trous — peut être déformée, tordue, découpée et recousue. Ces soixante et onze
pages, écrites à l'automne 1983, cherchent à faire l'inventaire complet de ces
manipulations et à comprendre comment elles s'enchaînent.

Le problème est le suivant. Pour chaque type de surface (un genre $g$, un
nombre $\nu$ de trous), il y a un groupe qui mesure toutes les façons de la
renvoyer sur elle-même, à déformation continue près : le \emph{groupe
modulaire}, ou groupe de Teichmüller. Ces groupes ne vivent pas isolés. Coudre
deux trous ensemble fait passer d'une surface à une autre ; découper le long
d'un cercle fait l'inverse ; boucher un trou par un disque aussi. L'ensemble de
ces surfaces et de ces opérations forme ce que Grothendieck appelle la
\emph{tour de Teichmüller}, et la question est d'en donner une description
purement combinatoire — avec des graphes, des ensembles finis, des
permutations — qui la reconstitue entièrement.

L'idée qui arrive est celle d'un jeu de construction. On découpe toute surface
en pièces élémentaires — des « pantalons », sphères à trois trous — et l'on
espère que tout se ramène à quelques pièces : celles de « dimension modulaire »
$0$ et $1$ pour engendrer, celles de dimension $2$ pour les relations. D'où les
haltes du dossier. Il étudie d'abord les deux courbes elliptiques qui ont plus
de symétries que les autres, l'une construite sur un carré, l'autre sur un
hexagone ; puis les configurations de quatre points sur une sphère, et la
correspondance entre le tore à un trou et la sphère à quatre trous ; puis, en
dimension $2$, les configurations de cinq points, dont il dresse un tableau de
$542$ structures « spéciales » reliées par $1620$ flèches, à la manière d'une
carte routière dont il faudrait trouver toutes les boucles. Vient ensuite la
partie la plus suivie, dix-sept pages qu'il numérote : les « groupoïdes de
Teichmüller », où les surfaces sont munies de cercles de découpage et de points
de repère, où chaque surface se résume en un graphe pondéré — un sommet par
morceau, une arête par cercle —, et où sont définies une à une les opérations
élémentaires : effacer des repères, en ajouter, effacer un cercle, boucher un
trou, découper, recoller. Le dossier finit sur un programme — un « petit
théorème de dévissage » qui dirait que tout se reconstitue à partir des
surfaces connexes — et sur des dessins de cartes découpées en pantalons.

On y voit une manière de travailler : avant tout théorème, nommer les objets
et les gestes (« gommage », « bouchage de trous », « surmarquage »), les
compter, les ranger, dire lesquels sont réversibles ; puis s'arrêter devant ce
qui serait trop gros — un modèle qui contiendrait « toutes les cordes » : « ça
fait un peu gros ! ».

Ces pages précèdent de quelques semaines l'\emph{Esquisse d'un programme}
(janvier 1984), où la tour de Teichmüller, le « jeu de Lego-Teichmüller » et le
principe des deux premiers étages sont exposés ; elles en montrent l'atelier.
Pour aller plus loin : groupe modulaire d'une surface, décomposition en
pantalons, complexe de Hatcher–Thurston, twist de Dehn, courbes elliptiques à
multiplication complexe, espaces de modules $\mathcal{M}_{0,4}$, $\mathcal{M}_{1,1}$,
$\mathcal{M}_{0,5}$, graphes stables, opérades modulaires, dessins d'enfants.

\keywords{Teichmüller tower, mapping class groupoid, Lego-Teichmüller game, pants decomposition, Hatcher–Thurston complex, boundary Dehn twist, fractional Dehn twist, CM elliptic curve, moduli space of five-pointed rational curves, real structure, stable graph, modular operad, Dehn–Thurston coordinates, cartographic group, dessins d'enfants}
\end{resume}

\pagerange{1}{71}
\section*{Le fil du dossier, et les conventions}

\subsection*{Les stations}

Le dossier n'a pas un argument unique : c'est un dossier de travail, sur des
papiers divers (feuillets blancs, dos de listings d'ordinateur, papier kraft),
dont deux pages seulement sont datées de sa main — la page 40, « Oct.\ 1983 »,
et la page 1, « 29.11.83 ». Mais tous ses morceaux servent un même projet, et
l'ordre des archivistes, qu'on suit, le laisse voir.

\begin{itemize}
\item \textbf{Groupes d'une surface à bord torique (pages 1 à 3)}, datées du
29 novembre 1983 : le groupe $SRT(X)$, sa filtration, les twists fractionnaires
$\mathbb{R}^{I}$, les marquages.
\item \textbf{Les deux courbes elliptiques exceptionnelles (pages 4 à 14, ses
pages 1 à 6)} : la courbe « carrée » et la courbe « hexagonale », leurs points
fixes, leurs décompositions cellulaires, leurs quotients par $\pm 1$.
\item \textbf{$\mathcal{M}_{0,4}$ et $\mathcal{M}_{1,1}$ (pages 16 à 20)} : une
carte sur la sphère, les quatre types de points-base et de chemins
élémentaires, et le diagramme reliant les groupes $\mathcal{T}_{0,4}$ et
$\mathcal{T}_{1,1} = SL_2(\mathbb{Z})$.
\item \textbf{Cartes, pièces et dimension modulaire (pages 22 à 27)} :
décomptes de cellules, les pièces rangées par dimension modulaire, les cartes
trivalentes et les revêtements de $\mathbb{P}^{1}$.
\item \textbf{$\mathcal{M}_{0,5}$ (pages 28 à 38)} : le tableau des neuf types
de structures spéciales sur cinq points, $542$ sommets et $1620$ flèches, un
programme de travail ; la classification des systèmes de cinq points à
symétrie (ses pages 1 à 3) ; le paramétrage du cas involutif (pages 37, 36, 38,
dans cet ordre).
\item \textbf{Le catalogue des opérations et la tour (pages 39 à 42)} : la
première liste des opérations (page 40, octobre 1983), une seconde (page 41),
et « les douze avatars principaux de la Tour de Teichmüller » (page 42).
\item \textbf{Groupoïdes de Teichmüller (pages 43 à 60, ses pages 1 à 17)} :
la partie suivie du dossier — définitions, foncteur vers les graphes,
graphes à repères, opérations élémentaires, récapitulation.
\item \textbf{Le programme de dévissage (pages 61 à 64, ses pages 18 à 21)}.
\item \textbf{Découpages en pantalons (pages 65 à 71)}, avec, page 70, une
autre mouture des définitions des pages 43 à 45.
\end{itemize}

Trois suites paginées par lui coexistent et ne se continuent pas : ses pages 1
à 6 (pages 4 à 14), ses pages 1 à 3 (pages 33 à 35), ses pages 1 à 21 (pages 43
à 64). La page 70 recommence les définitions de la page 43 avec d'autres
lettres ; c'est une autre mouture et on la lit à part, sans la fondre dans la
première. Rien ne dit laquelle est la plus ancienne.

\subsection*{Les conventions}

\emph{Groupoïdes.} Le dossier écrit $\mathcal{T}$ pour le groupoïde orienté
(flèches conservant l'orientation) et $\mathcal{T}^{\pm}$ pour le biorienté ;
une note de la page 43 propose d'inverser cette convention, sans que la suite
la suive. On écrit partout l'exposant : $\mathcal{T}^{+}$ pour l'orienté,
$\mathcal{T}^{\pm}$ pour le biorienté\footnote{Sauf dans les noms génériques
(« la tour $\mathcal{T}$ », les losanges des pages 40 et 41), où la variance
ne compte pas.}. Le préfixe $S_{*}$ (« spécial ») indique des points de
rigidification, le suffixe $D$ (« divisé ») des cercles de découpage. Un
groupoïde à un seul objet de groupe $G$ est noté $[G]$ ; en particulier
$[\pm 1]$.

\emph{Graphes.} Dans un graphe, $S$ est l'ensemble des sommets, $\vec{A}$
celui des demi-arêtes, $\sigma$ l'involution, $o : \vec{A} \to S$ l'origine ;
$A = \vec{A}/\sigma = A_0 \sqcup A_\ell$ se partage en arêtes liées et arêtes
libres (les points fixes de $\sigma$). Le poids des sommets est le genre
$\gamma : S \to \mathbb{N}$, celui des arêtes le nombre de repères
$r : A \to \mathbb{N}$. La lettre $S$ étant prise par les sommets, l'ensemble
fini de points marqués est $R$ partout ; la page 70 l'appelle $S$ et on la
traduit.

\emph{Deux usages de « $T$ ».} Aux pages 1 à 3, $T(X)$, $SRT(X)$, $SUT(X)$ sont
des \emph{groupes} attachés à une surface $X$, et $T_i$ est un cercle du bord ;
à partir de la page 19, $\mathcal{T}$ désigne des groupes ou groupoïdes de
Teichmüller. On garde les deux, la typographie les distingue.

\emph{Opérations.} On numérote les opérations élémentaires comme la
récapitulation de la page 60 : $0^\circ$ somme, $1^\circ$ à $3^\circ$
gommage, gommage partiel, surmarquage, $4^\circ$ gommage d'un cercle de
découpage, $5^\circ$ bouchage d'un trou, $6^\circ$ découpage, $7^\circ$
recollement.

\emph{Le mot « dimension modulaire ».} Pour une surface de genre $g$ à $\nu$
trous, c'est $3g - 3 + \nu$, la dimension complexe de l'espace des modules
$\mathcal{M}_{g,\nu}$ ; la condition « anabélienne » des pages 40 et 41 est
$2g - 2 + \nu > 0$, la surface est hyperbolique.

\pagerange{1}{3}
\section*{Groupes d'une surface à bord torique (pages 1 à 3)}

\pagerange{1}{1}
\subsection*{Le groupe $SRT(X)$ et sa filtration}

Soit $X$ une surface compacte orientée à bord, et supposons chaque composante
connexe $T_i$ du bord ($i \in I = \pi_0(\partial X)$) munie d'une
\emph{structure torique} : une structure de torseur sous le groupe du cercle
$\mathbb{U} = \{ z \in \mathbb{C} \mid |z| = 1 \}$, autrement dit une
paramétrisation du cercle définie à rotation près. On pose
\[
SRT(X) = \mathrm{Aut}(X)/\mathrm{Aut}^{!}(X)_0 ,
\]
où $\mathrm{Aut}(X)$ est le groupe des homéomorphismes de $X$ qui respectent
l'orientation et les structures toriques\footnote{« respectant » est une
lecture incertaine. La page écrit $\mathrm{Aut}^{!}(\vec{X})$, la flèche
marquant l'orientation.}, et $\mathrm{Aut}^{!}(X)_0$ la composante neutre du
sous-groupe $\mathrm{Aut}^{!}(X)$ de ceux qui induisent l'identité sur
$\partial X$. Le quotient
\[
S_{*}T^{!}(X) = \mathrm{Aut}^{!}(X)/\mathrm{Aut}^{!}(X)_0
\]
est ce qu'on appelle aujourd'hui le groupe modulaire de $X$ \emph{relatif au
bord}, $\mathrm{Mod}(X, \partial X)$ : homéomorphismes égaux à l'identité sur
le bord, à isotopie fixant le bord près.

L'action sur le bord donne une suite exacte
\[
1 \to S_{*}T^{!}(X) \to SRT(X) \to \mathrm{Aut}\bigl((T_i)_{i \in I}\bigr),
\]
et le groupe des automorphismes de la famille de torseurs
$(T_i)$ est une extension de $\mathfrak{S}_I$ par $\mathbb{U}^{I}$. D'où une
filtration de $SRT(X)$ à trois étages, de quotients successifs
$S_{*}T^{!}(X)$, $\mathbb{U}^{I}$ et $\mathfrak{S}_I$\footnote{Le dernier
quotient est $\mathfrak{S}_I$ entier quand $X$ est connexe : toute permutation
des trous et toute rotation des bords se réalisent alors par un homéomorphisme.
Pour $X$ non connexe, il faut le remplacer par le sous-groupe des permutations
qui respectent les composantes et leur type.}.

\emph{Les twists de bord.} Sur un collier $T_i \times [0,1]$ de chaque cercle du
bord, on peut tourner le cercle $T_i \times \{t\}$ d'un angle $2\pi\theta_i t$ ;
pour $\theta = (\theta_i) \in \mathbb{R}^{I}$ cela définit un élément de
$SRT(X)$ qui tourne le bord $T_i$ de $\theta_i$ tours. On obtient ainsi un
sous-groupe canonique $\mathbb{R}^{I} \subset SRT(X)$ — des \emph{twists de
Dehn fractionnaires} —, dont la partie entière $\mathbb{Z}^{I}$ est formée des
twists de Dehn le long des cercles du bord, qui sont dans $S_{*}T^{!}(X)$.
Si chaque composante de $X$ est hyperbolique ($2g - 2 + \nu > 0$),
\[
\mathbb{R}^{I} \cap S_{*}T^{!}(X) = \mathbb{Z}^{I}, \qquad
SRT^{!}(X) = \mathbb{R}^{I} \cdot S_{*}T^{!}(X),
\]
où $SRT^{!}(X)$ est le noyau de $SRT(X) \to \mathfrak{S}_I$, et
$\mathbb{R}^{I}$ est central dans $SRT^{!}(X)$, puisqu'un élément de
$S_{*}T^{!}(X)$ peut être choisi égal à l'identité sur les colliers. Ainsi
$SRT^{!}(X)$ est le produit central
\[
SRT^{!}(X) \simeq \bigl(\mathbb{R}^{I} \times S_{*}T^{!}(X)\bigr)/\mathbb{Z}^{I},
\]
$\mathbb{Z}^{I}$ plongé antidiagonalement\footnote{L'hypothèse est ajoutée. Sur
un disque, le twist de bord est trivial (lemme d'Alexander) ; sur un anneau, les
twists le long des deux bords sont égaux ; dans ces cas $\mathbb{Z}^{I} \to
S_{*}T^{!}(X)$ n'est pas injectif. La condition est celle que la page 40 appelle
« anabélienne ».}. Le treillis que la page dessine est donc :

\begin{tikzcd}[column sep=small, row sep=small]
 & SRT(X) \arrow[d, no head, "\mathfrak{S}_I"] & \\
 & SRT^{!}(X) \arrow[dl, no head, "T^{!}(X)"'] \arrow[dr, no head, "\mathbb{U}^{I}"] & \\
\mathbb{R}^{I} \arrow[dr, no head, "\mathbb{U}^{I}"'] & & S_{*}T^{!}(X) \arrow[dl, no head, "T^{!}(X)"] \\
 & \mathbb{Z}^{I} &
\end{tikzcd}

où les étiquettes sont les quotients, et $T^{!}(X) = S_{*}T^{!}(X)/\mathbb{Z}^{I}$
est le groupe modulaire pur de $X$ à isotopie libre sur le bord — celui de la
surface où l'on aurait remplacé chaque trou par une piqûre.

\emph{L'extension $(*)$.} Posons $SUT(X) = SRT(X)/\mathbb{Z}^{I}$ et $T(X) =
SRT(X)/\mathbb{R}^{I}$, le groupe modulaire de $X$ où les bords peuvent tourner
et être permutés. Comme $\mathbb{R}^{I}/\mathbb{Z}^{I} = \mathbb{U}^{I}$, on a
l'extension
\[
(*)\qquad 1 \to \mathbb{U}^{I} \to SUT(X) \to T(X) \to 1,
\]
et, comme le note la marge, $SUT(X) \simeq T(X) \times_{\mathfrak{S}_I}
\mathrm{Aut}\bigl((T_i)_{i\in I}\bigr)$ : l'application vers ce produit fibré
est injective, car son noyau est contenu dans $\mathbb{U}^{I}$ qui s'envoie
isomorphiquement sur les rotations, et surjective, car on corrige un relèvement
par une rotation\footnote{Une note oblique de la marge gauche, presque
entièrement illisible, revient sur le rôle de $S_{*}T^{!}(X)$ dans la définition
de $T(X)$ ; on n'en tire rien.}.

\pagerange{1}{2}
\subsection*{Marquages et scindages}

Un \emph{marquage} du bord est un point $\underline{s} = (s_i) \in \prod_{i \in
I} T_i$, un point sur chaque cercle. Il définit un scindage
$\varphi_{\underline{s}}$ de $(*)$ : au-dessus d'un élément de $T(X)$ de
permutation $\pi$, on prend l'élément de $SUT(X)$ dont l'action sur le bord
envoie chaque $s_i$ sur $s_{\pi(i)}$. Le groupe $S_{*}T(X, \underline{s})$ des
classes d'homéomorphismes qui respectent l'ensemble des $s_i$ (à isotopie
fixant les $s_i$ près) s'identifie à l'image inverse de
$\varphi_{\underline{s}}(T(X))$ par $SRT(X) \to SUT(X)$\footnote{Le premier
groupe est écrit « $SUR(X)$ » ; lu $SRT(X)$. L'identification utilise que le
groupe des homéomorphismes du cercle fixant un point est contractile, de sorte
qu'isotopie fixant $s_i$ et isotopie fixant $T_i$ donnent le même quotient.}.
On a alors
\[
SRT(X) = \mathbb{R}^{I} \cdot S_{*}T(X, \underline{s}), \qquad
\mathbb{R}^{I} \cap S_{*}T(X, \underline{s}) = \mathbb{Z}^{I},
\]
avec $\mathbb{R}^{I}$ distingué ; sur les parties pures, c'est le produit
central $\bigl(\mathbb{R}^{I} \times S_{*}T^{!}(X, \underline{s})\bigr)/
\mathbb{Z}^{I}$\footnote{La page écrit $S_{*}T^{!}(X, \underline{s})
\wedge^{\mathbb{Z}^{I}} \mathbb{R}^{I} \simeq SRT(X)$, l'exposant $!$ étant un
signe surchargé lu sans certitude. Avec l'exposant, le membre de gauche est la
partie pure $SRT^{!}(X)$ ; sans lui, le « produit contracté » n'est qu'un produit
de sous-groupes, $\mathbb{R}^{I}$ n'étant plus central : les permutations le
normalisent sans le centraliser.}.

Changer $\underline{s}$ en $\underline{s}' = u \cdot \underline{s}$, $u \in
\mathbb{U}^{I}$, change $\varphi_{\underline{s}}$ en $\mathrm{int}(u) \circ
\varphi_{\underline{s}}$, donc, si $\tilde u \in \mathbb{R}^{I}$ relève $u$,
\[
S_{*}T(X, \underline{s}') = \mathrm{int}(\tilde u)\, S_{*}T(X, \underline{s}).
\]
Ainsi $S_{*}T(X, \underline{s})$ dépend d'un choix non canonique, tandis que
$SRT(X)$ ne dépend que des structures toriques. La page voulait ensuite
réaliser le torseur $\prod T_i$ comme un ensemble de scindages de $(*)$ ; le
passage est en grande partie illisible, et l'application $\underline{s} \mapsto
\varphi_{\underline{s}}$, équivariante, n'est pas injective en général : pour $X$
connexe, tourner tous les $s_i$ du même angle ne change pas
$\varphi_{\underline{s}}$. On n'affirme donc rien de plus\footnote{Cette
non-injectivité est notre remarque ; elle vient de ce qu'une rotation diagonale
commute aux permutations.}.

\pagerange{3}{3}
\subsection*{Multimarquages}

Un \emph{multimarquage} est la donnée, pour chaque $i$, d'un entier $n_i \geq 1$
et d'un sous-torseur $T_i(n_i) \subset T_i$ sous $\mu_{n_i} \subset \mathbb{U}$
— $n_i$ points régulièrement espacés —, c'est-à-dire d'un point $s_i$ de
$T_i^{\wedge n_i} = T_i/\mu_{n_i}$. Le groupe de Teichmüller multimarqué se
plonge dans $SRT(X)$ : on choisit un point $s^0_i$ dans chaque $T_i(n_i)$ et
\[
S_{*}T(X, \underline{s}) = S_{*}T(X, \underline{s}^{0}) \cdot \prod_{i \in I}
\tfrac{1}{n_i}\mathbb{Z}_i \ \subset SRT(X),
\]
où $\mathbb{Z}_i$ est le $i$-ème facteur de $\mathbb{Z}^{I} \subset
\mathbb{R}^{I}$ : les twists d'un $n_i$-ième de tour permutent circulairement
les points de $T_i(n_i)$\footnote{Pour que ce produit soit un groupe, il faut ne
garder dans $S_{*}T(X, \underline{s}^{0})$ que les éléments dont la permutation
respecte la fonction $i \mapsto n_i$ ; la page ne le dit pas.}. Ces
twists fractionnaires reviennent aux pages 40 et 56, sous le nom de twist
relatif et de surmarquage uniforme.

\pagerange{4}{14}
\section*{Les deux courbes elliptiques exceptionnelles (pages 4 à 14)}

Les courbes elliptiques complexes ont toutes l'automorphisme $x \mapsto -x$ ;
deux seulement, à isomorphisme près, en ont d'autres : celle de module $j =
1728$, $\mathbb{C}/\mathbb{Z}[i]$, et celle de module $j = 0$,
$\mathbb{C}/\mathbb{Z}[\zeta]$ avec $\zeta = e^{2i\pi/3}$. Ce sont les deux
points orbifolds de $\mathcal{M}_{1,1}$ ; les pages 4 à 14 (ses pages 1 à 6) les
décrivent de façon combinatoire\footnote{Ce sont les courbes à multiplication
complexe par les entiers de $\mathbb{Q}(i)$ et de $\mathbb{Q}(\sqrt{-3})$. Le
dossier 141, page 9, dessine la droite des modules avec ses deux points
exceptionnels ; voir sa lecture modernisée.}.

\pagerange{4}{6}
\subsection*{La courbe carrée}

La catégorie des courbes elliptiques (pointées) isomorphes à
$\mathbb{E}\mathrm{l}_0 = \mathbb{C}/\mathbb{Z}[i]$, avec leurs isomorphismes,
est un groupoïde connexe de groupe $\mathrm{Aut}(\mathbb{E}\mathrm{l}_0) = \mu_4
= \mathbb{Z}[i]^{*}$ ; elle est équivalente à la catégorie des carrés
combinatoires orientés, c'est-à-dire des torseurs $S$ sous $\mathbb{Z}/4 \simeq
\mu_4$, par
\[
S \longmapsto \mathbb{E}\mathrm{l}_0(S) = \mathbb{E}\mathrm{l}_0 \wedge^{\mu_4}
S .
\]
Concrètement : on réalise le carré $S$ comme carré euclidien de centre $o$ dans
un plan $V_S$ ; ses sommets engendrent un réseau $\Gamma_S$, ses diagonales
donnent deux axes de coordonnées définis aux signes près, la structure
euclidienne orientée de $V_S$ une structure complexe, et $X_S =
\mathbb{E}\mathrm{l}_0(S) = V_S/\Gamma_S$.

\emph{Points d'ordre $2$.} Le générateur $i$ de $\mu_4 = \mathrm{Aut}(X_S)$ fixe
exactement deux points,
\[
X_S^{\,i} = \{ x \mid ix = x \} = \{ e, \eta_S \},
\]
car $(1 - i)x \in \mathbb{Z}[i]$ définit un sous-groupe d'ordre $|1-i|^2 = 2$ ;
$\eta_S$ est l'unique point d'ordre $2$ fixé par $\mathrm{Aut}(X_S)$ (dans
$\mathbb{C}/\mathbb{Z}[i]$, c'est $(1+i)/2$), et $i$ échange les deux autres,
$1/2$ et $i/2$.

\emph{Cellules.} $X_S$ a une décomposition cellulaire à $1$ sommet $e$, $2$
arêtes (les deux mailles du réseau) et $1$ face. Son image inverse par la
multiplication par $2$, de degré $4$, a $4$ sommets (les points d'ordre $2$),
$8$ arêtes et $4$ faces. Les deux sont invariantes par $\mu_4$. La seconde se
prête au quotient par $\pm 1$ : $X_S/\pm 1$ est une sphère — l'« oreiller »,
revêtu doublement par le tore et ramifié aux quatre points d'ordre $2$ —,
formée de deux carrés recollés le long d'un équateur que les images de $e$,
$\eta$, $\eta'$, $\eta''$ divisent en quatre arêtes\footnote{L'exposant
qu'il met aux images ($e^{q}$, \ldots) est lu $q$ sans certitude ; le mot
au-dessus de « carrés » est illisible.}.

\pagerange{8}{14}
\subsection*{La courbe hexagonale}

De même, les courbes isomorphes à $\mathbb{E}\mathrm{l}_1 =
\mathbb{C}/\mathbb{Z}[\zeta]$ forment un groupoïde équivalent à celui des
hexagones orientés, torseurs $S$ sous $\mathbb{Z}/6 \simeq \mu_6 =
\mathbb{Z}[\zeta]^{*}$, par $S \mapsto \mathbb{E}\mathrm{l}_1 \wedge^{\mu_6} S$.
On réalise l'hexagone $S$ dans un plan euclidien orienté $V_S$, d'origine son
barycentre ; ses sommets engendrent un réseau $\Gamma_S$ et $X_S = V_S/\Gamma_S$
\footnote{La page appelle d'abord $V_S$ le réseau engendré par $S$, puis écrit
$V_S/\Gamma_S$ ; c'est $\Gamma_S$.}.

\emph{Points d'ordre $2$.} $\mu_6$ opère sur les trois points d'ordre $2$ non
nuls, $-1$ trivialement ; ils forment un torseur sous $\mu_6/\pm 1 \simeq
\mathbb{Z}/3$, identifié à $S/\pm 1$ en envoyant un sommet $s$ sur l'image du
milieu de $[0, s]$\footnote{La page écrit $\zeta = \exp\frac{2\pi}{6}$ pour un
générateur ; il faut $\exp\frac{2i\pi}{6}$. La page 8 a bien $\zeta = \exp
2i\pi/3$.}.

\emph{Points fixes.} Les points fixes de $\zeta$ (d'ordre $3$) sont les $x$ avec
$(1 - \zeta)x \in \mathbb{Z}[\zeta]$, et $|1 - \zeta|^2 = 3$ : il y en a trois,
l'origine $e$ et deux points d'ordre $3$, symétriques, comme $\pm(1 - \zeta)/3$.
Dans l'hexagone, ce sont les centres des six triangles de la subdivision
centrale, qui se répartissent en deux classes de trois, congrus modulo $\Gamma_S$
\footnote{« régiments » est une lecture incertaine.} ; ils correspondent aux
deux triangles inscrits dans l'hexagone (sommets deux à deux non adjacents).
Sous $\mu_6$ tout entier, seul $e$ est fixe, car $1 + \zeta^2 = -\zeta$ est une
unité.

\emph{Cellules.} Décomposition canonique : un sommet, trois arêtes (les trois
diagonales de l'hexagone, ou les trois points d'ordre $2$ non nuls), deux faces
triangulaires, contenant chacune un des deux points fixes d'ordre $3$. Doublée
par la multiplication par $2$ : $4$ sommets, $12$ arêtes, $8$ faces. Son
quotient par $\pm 1$ a $4$ sommets, $6$ arêtes et $4$ faces triangulaires :
c'est un tétraèdre, orienté, avec un sommet privilégié (l'image de $e$), et son
groupe d'automorphismes est $\mathbb{Z}/3 = \mu_6/\pm 1$, ce que la page vérifie
(« OK ! »).

Les deux quotients — la sphère à deux carrés et le tétraèdre — sont les
structures « carrée » et « tétraédrale » sur quatre points d'une sphère qui
reviennent aux pages 17 à 30.

\pagerange{16}{20}
\section*{$\mathcal{M}_{0,4}$ et $\mathcal{M}_{1,1}$ (pages 16 à 20)}

La page 16, une couverture, porte seulement « $M_{0\,4}$ » et « $M_{1,1}$ » :
l'espace des modules des sphères à quatre points et celui des courbes
elliptiques, les deux pièces de dimension modulaire $1$. Les pages qui suivent
sont des feuilles d'esquisses, aux encres de couleur et au crayon.

\pagerange{17}{18}
\subsection*{Une carte, des points-base, des chemins}

La page 17 compte les cellules d'une carte sur la sphère, dessinée à partir
d'un tétraèdre tronqué et de sa subdivision barycentrique : $4 + 6 + 12 = 22$
sommets (quatre et six d'ordre $6$, douze d'ordre $3$), $24$ triangles et $4$
hexagones, soit $28$ faces, et donc $48$ arêtes\footnote{La page écrit le
nombre d'arêtes sur une surcharge (« $4.9$ », « $48$ » ?). $48$ est imposé
deux fois : par la formule d'Euler, $22 - 48 + 28 = 2$, et par le décompte
$(24 \cdot 3 + 4 \cdot 6)/2 = 48$ ; la somme des valences, $4\cdot 6 + 6\cdot 6
+ 12 \cdot 3 = 96$, le confirme.}.

Il y liste ensuite quatre types de « points-base » — des configurations
spéciales de quatre points $I$ sur la sphère : trois structures carrées, deux
structures tétraédrales (correspondant aux deux orientations), six structures
« bipantalonniques » et six structures « pantalonniques » de genre
$0$\footnote{« de genre $0$ » est une lecture incertaine (« de genre gois »).
On peut reconnaître dans les deux premiers types les trois points harmoniques
($\lambda = -1, 1/2, 2$) et les deux points équianharmoniques ($\lambda = -\zeta,
-\zeta^2$) de l'espace des modules de quatre points numérotés ; les douze autres
semblent attachés aux trois pointes, deux fois deux par pointe. Ces
identifications sont les nôtres.} —, puis quatre types de chemins élémentaires
entre eux : quart de twist ($2 \times 3 \times 4$ chemins), sixième de tour sur
l'équateur, pôle vers trou voisin, pôle vers équateur voisin ($2 \times 2 \times
3$ chacun) ; deux types composés, trou vers trou (« tiers de tour », $2 \times
3$) et demi-twist ; et le twist complet. Avec les deux sens de parcours, cela
fait huit types de déformations. Les « relations » sont annoncées et non
écrites.

La page 18 ajoute que, pour passer aux structures pantalonniques, il faut la
subdivision barycentrique complète de la carte carrée ; que dans les deux cas
non dégénérés (carré, tétraèdre) toutes les positions adjacentes sont dans la
subdivision barycentrique ; et que la subdivision barycentrique du tétraèdre
contient celles des trois cartes carrées associées.

\pagerange{19}{20}
\subsection*{Le tore à un trou et la sphère à quatre trous}

Soit $I = \{a, b, c, d\}$. Ses trois partitions en deux paires,
\[
Q_0 = \{ad, bc\}, \qquad Q_1 = \{bd, ac\}, \qquad Q_\infty = \{cd, ab\},
\]
forment un ensemble $J$ à trois éléments — les trois pointes de
$\mathcal{M}_{0,4}$, où deux points se rejoignent deux à deux —, d'où la suite
exacte
\[
1 \to V(J) \to \mathfrak{S}_I \to \mathfrak{S}_J \to 1,
\]
de noyau le groupe de Klein. Les orientations de $I$ et de $J$ se correspondent
($\omega_I \simeq \omega_J$), car $\mathfrak{S}_4 \to \mathfrak{S}_3$ respecte
la signature\footnote{La page écrit aussi des isomorphismes entre ensembles à
deux éléments, $\omega_I \simeq \omega_J \simeq C_Q \simeq J \setminus \{Q\}$,
et une correspondance (tétraèdre orienté $+$ diagonale) $\leftrightarrow$ (carré
$+$ codiagonale), qu'on ne détaille pas.}.

Le cadre $\rho^3 = \sigma^2 = 1$, $\rho\sigma = \ell$ est la présentation de
$PSL_2(\mathbb{Z}) \simeq \mathbb{Z}/3 * \mathbb{Z}/2$, où $\ell$ est la
transvection parabolique, qu'on retrouve au dossier 145 comme $\mathcal{E}^{+}
= \langle \sigma_0, \rho \mid \sigma_0^2 = \rho^3 = 1 \rangle$. Le diagramme
du bas de la page relie les groupes modulaires de la sphère à quatre trous et
du tore à un trou ; ce qui est vrai est
\[
\begin{array}{ccccccccc}
1 \to & \mathcal{T}^{!}_{0,4} & \to & \mathcal{T}_{0,4} & \to & \mathfrak{S}_4 & \to 1 \\
 & \downarrow{\scriptstyle\wr} & & \downarrow & & \downarrow & \\
1 \to & \pi_{0,3} & \to & \mathcal{T}'_{0,4} & \to & \mathfrak{S}_3 & \to 1 \\
 & \uparrow{\scriptstyle\wr} & & \uparrow & & \uparrow & \\
1 \to & N & \to & \mathcal{T}_{1,1} & \to & \widetilde{\mathfrak{S}}_3 & \to 1
\end{array}
\]
où $\mathcal{T}_{0,4}$ est le groupe modulaire de la sphère à quatre points,
$\mathcal{T}^{!}_{0,4}$ son sous-groupe pur, libre de rang $2$, $\mathcal{T}'_{0,4}
= \mathcal{T}_{0,4}/V \simeq PSL_2(\mathbb{Z})$ son quotient par le sous-groupe de
Klein qui relève $V(J)$, $\mathcal{T}_{1,1} = SL_2(\mathbb{Z})$ celui du tore à
un trou, et la flèche $\mathcal{T}_{1,1} \to \mathcal{T}'_{0,4}$ le quotient par
l'involution hyperelliptique $-1$. La ligne du bas demande un mot. Écrivons
$SL_2(\mathbb{Z}) = \langle S, T \mid S^4 = 1,\ S^2 = (ST)^3 \rangle$, et soit
$\widetilde{\mathfrak{S}}_3 = \langle x, a \mid x^4 = 1,\ x^2 = a^3,\ x a x^{-1}
= a^{-1} \rangle$ le groupe dicyclique d'ordre $12$, extension centrale de
$\mathfrak{S}_3$ par $\mathbb{Z}/2$. Le morphisme $S \mapsto x$, $ST \mapsto a$
est surjectif, envoie $-1 = S^2$ sur l'élément central $x^2 \neq 1$, et il est
injectif sur les sous-groupes finis $\langle S \rangle$ et $\langle ST \rangle$ ;
son noyau $N$ est donc sans torsion, d'indice $12$, libre de rang $2$, et ne
contient pas $-1$. Comme $\mathfrak{S}_3$ n'est quotient de $PSL_2(\mathbb{Z})$
que d'une seule façon, $N \cdot \{\pm 1\} = \Gamma(2)$, et $N$ s'envoie
isomorphiquement sur $\pi_{0,3} = \Gamma(2)/\pm 1$ : la ligne du bas de la page
est exacte, avec son « $\wr$ »\footnote{Le premier terme est lu
$\pi_{0,3}(2, \mathbf{Z})$ sans certitude ; c'est $N$, un relèvement de
$\pi_{0,3}$ dans $SL_2(\mathbb{Z})$, distinct du sous-groupe de congruence
$\Gamma(2) = N \times \{\pm 1\}$. L'annotation $\widetilde{\mathfrak{S}}_3
\simeq SL(2, \mathbb{Z}/4)/\pm$ ne convient pas : ce groupe est d'ordre $24$
(c'est $\mathfrak{S}_4$, vérifié par les ordres des éléments), non $12$.
L'argument par le groupe dicyclique est le nôtre, ainsi que l'identification de
$\mathcal{T}'_{0,4}$ comme quotient de $\mathcal{T}_{0,4}$ par $V$, que suggère la
ligne du milieu.}. Les annotations
$\mathfrak{S}_4 \simeq \mathrm{Aff}(2, \mathbb{F}_2)$ et $\mathfrak{S}_3 \simeq
SL(2, \mathbb{F}_2)$ sont exactes, et la première s'explique par les pages 4 à
6 : les quatre points de ramification de $E \to E/\pm 1 = \mathbb{P}^{1}$ sont
les images des points d'ordre $2$ de $E$, qui forment un plan affine sur
$\mathbb{F}_2$\footnote{Pour la suite exacte $1 \to \Gamma(2)/\pm 1 \to
PSL_2(\mathbb{Z}) \to \mathfrak{S}_3 \to 1$, voir aussi la lecture modernisée du
dossier 141, page 9.}.

La page 20 ne porte qu'un dessin, un tétraèdre et un octaèdre traversé d'un
axe.

\pagerange{22}{27}
\section*{Cartes, pièces et dimension modulaire (pages 22 à 27)}

\pagerange{22}{23}
\subsection*{Décomptes et cartes canoniques}

La page 22 reprend les tores des pages 4 à 14 : un tétraèdre à huit triangles
(quatre sommets d'ordre $6$, douze arêtes) ; les groupes $\mu_6 \times
(\mathbb{Z}/2)^2$ et $\mu_4 \times (\mathbb{Z}/2)^2$ — les automorphismes des
tores hexagonal et carré joints aux translations par les points d'ordre $2$ ;
le carré orienté, quatre carrés, huit arêtes, quatre sommets ; et un carré de
correspondances entre $E_{\varepsilon,\pm}$, $E_Q$ et $E_H$, laissé sans texte.

La page 23 note $\mathrm{diag}_Q \simeq \omega_Q$ et compte les « cartes
canoniques » de type $(0,4)$ sur un ensemble $I$ de quatre points : deux
tétraédrales, trois carrées, six « assemblages de $J$ droits » et six
« gauches »\footnote{La somme de la liste est $2 + 3 + 6 + 6 = 17$, qui est le
nombre des points-base de la page 17 ; la page écrit « $14$ » sur un « $17$ »,
avec « $14 = 3 \times 4 + 2$ ». On ne tranche pas.}.

\pagerange{24}{24}
\subsection*{Les pièces, par dimension modulaire}

La page 24 est un tableau, qu'on complète par la dimension modulaire $d = 3g -
3 + \nu$ qu'il note en accolade :
\[
\begin{array}{llll}
d = 0 & (0,3) & \text{triangle } A & \Sigma_A \\
d = 1 & (0,4) & \text{tétraèdre orienté } \vec{T},\ \text{carré } Q & \Sigma_T,\ \Sigma_Q \\
d = 1 & (1,1) & \text{carré orienté},\ \text{hexagone orienté } \vec{H} & E_{\vec{Q}},\ E_{\vec{H}} \\
d = 2 & (0,5),\ (1,2) & &
\end{array}
\]
Une pièce combinatoire — triangle, tétraèdre, carré, hexagone — détermine une
sphère à trois ou quatre points spéciale, ou l'une des deux courbes
elliptiques exceptionnelles. C'est l'échafaudage de ce que l'\emph{Esquisse
d'un programme} appellera le jeu de « Lego-Teichmüller » : la tour engendrée
par les pièces de dimension modulaire $\leq 1$, ses relations venant de la
dimension $2$, c'est-à-dire de $\mathcal{M}_{0,5}$ et de
$\mathcal{M}_{1,2}$\footnote{L'\emph{Esquisse} est de janvier 1984 ; ces pages
sont de l'automne 1983. La quatrième lettre de la première colonne, un $Q$
fléché sur une surcharge, est transcrite faute de mieux ; « orienté » est biffé
après « carré » à la troisième ligne.}.

\pagerange{25}{27}
\subsection*{Cartes trivalentes et revêtements}

La page 25 énonce une correspondance entre cartes connexes dont tous les sommets
sont d'ordre $3$ ou $1$, cartes triangulées (par dualité) et revêtements de
$\mathbb{P}^{1}$ ramifiés seulement en $0, 1, \infty$, et écrit les relations
\[
\sigma_0^2 = \sigma_1^2 = \sigma_2^2 = (\sigma_0\sigma_2)^2 = (\sigma_0\sigma_1)^3 = 1 .
\]
Ce sont les relations du groupe cartographique augmentées de $(\sigma_0\sigma_1)^3
= 1$ : le groupe qu'elles présentent est le groupe modulaire étendu
$PGL_2(\mathbb{Z})$, dont les ensembles finis à action sont les cartes
trivalentes (non orientées), et son sous-groupe d'indice $2$ des mots pairs,
$PSL_2(\mathbb{Z})$, celui des cartes trivalentes orientées\footnote{Voir les
lectures modernisées du dossier 141, pages 9 et 10 (où ce groupe est noté
$\Gamma_{\infty,3}$), et du dossier 145, page 22 (« groupe cartographique de
Malgoire »).}. La correspondance avec les revêtements doit être précisée : les
cartes trivalentes orientées correspondent aux revêtements finis de
$\mathbb{P}^{1}$ non ramifiés hors de $0, 1, \infty$ dont les indices de
ramification divisent $3$ au-dessus d'un des points et $2$ au-dessus d'un autre
— les revêtements de la droite modulaire ; un revêtement quelconque ramifié en
trois points correspond à une carte bicolore (un dessin d'enfant), et devient
trivalent après composition avec l'application de degré $6$, $\lambda \mapsto
j(\lambda)$\footnote{La page écrit « $\simeq$ revêtements de $\mathbb{P}^{1}$
ramifiés seulement en $0, 1, \infty$ », le mot « revêtements » surchargé ou
barré. La restriction est de cette édition. Le théorème de Belyi (1979), qui
rend ces revêtements exactement ceux des courbes définies sur
$\overline{\mathbb{Q}}$, était paru.}. Les cartes non orientées, avec les
$\sigma_i$ sans point fixe, conduisent aux $\mathbb{R}$-courbes algébriques que
la page nomme.

Les pages 26 et 27 comptent des cellules. Deux paires de décomptes,
\[
(10, 16, 4) \Longrightarrow (16, 32, 8), \qquad (6, 8, 4) \Longrightarrow (8,
16, 8)
\]
(sommets, arêtes, faces hexagonales puis carrées), sont celles d'un revêtement
double ramifié en quatre sommets : $V' = 2V - 4$, $E' = 2E$, $F' = 2F$ ; le
second va de la sphère ($\chi = 2$) au tore ($\chi = 0$) — l'oreiller et le tore
carré des pages 4 à 6, à une subdivision près\footnote{La lecture par
Riemann–Hurwitz est la nôtre. La page note « dont 4 ramifiés d'ordre \ldots »
sans chiffre.}. Les autres décomptes, ($2$ hexagones, $9$ arêtes, $6$ sommets)
et ($4$ sommets, $6$ arêtes, $2$ hexagones), sont isolés.

\pagerange{28}{38}
\section*{$\mathcal{M}_{0,5}$ (pages 28 à 38)}

Cinq points sur une sphère : c'est la première situation de dimension modulaire
$2$, celle d'où doivent venir les relations.

\pagerange{30}{32}
\subsection*{Le tableau des structures spéciales (pages 30 à 32)}

Soit $K$ un ensemble de cinq points. La page 30 est un tableau de neuf types de
structures « spéciales » sur $K$, rangées sur trois niveaux qui sont ceux des
strates de $\overline{\mathcal{M}}_{0,5}$ : au niveau $2$, la sphère n'est pas
découpée ; au niveau $1$, un cercle la coupe en une pièce $(0,3)$ et une pièce
$(0,4)$ ; au niveau $0$, deux cercles la coupent en trois pantalons. Pour
chacune, $G$ est le groupe de ses automorphismes holomorphes et
$\widetilde{G}$ celui où l'on admet aussi les antiholomorphes ; le nombre de
structures d'un type est $120/|G|$.
\[
\begin{array}{lllll}
\text{niveau} & \text{type} & \text{nom} & G & \text{nombre} \\
2 & \mathrm{I}_{10} & \text{pentagonale} & \mathbb{D}_5 & 12 \\
2 & \mathrm{I}_{6} & \text{bitétraédrale} & \mathfrak{S}_3 & 20 \\
2 & \mathrm{I}_{4} & \text{pyramidale (carré orienté)} & \mathbb{Z}/4 & 30 \\
1 & \mathrm{II}_{2} & \text{carré-triangle droit} & \mathbb{Z}/2 & 60 \\
1 & \mathrm{II}'_{2} & \text{carré-triangle gauche} & \mathbb{Z}/2 & 60 \\
1 & \mathrm{II}_{1} & \text{tétraèdre-triangle} & 1 & 120 \\
0 & \mathrm{III}_{2} & \text{triple triangle droit} & \mathbb{Z}/2 & 60 \\
0 & \mathrm{III}'_{1} & \text{triple triangle semi-gauche} & 1 & 120 \\
0 & \mathrm{III}'_{2} & \text{triple triangle gauche} & \mathbb{Z}/2 & 60
\end{array}
\]
soit $62 + 240 + 240 = 542$ sommets. Dans chaque cas $\widetilde{G} = G \times
\{1, \sigma\}$, sauf pour $\mathrm{I}_4$ où $\widetilde{G} = G$ : « sauf
$\mathrm{I}_4$, toutes ces structures sont réelles »\footnote{Il s'agit des
automorphismes de la \emph{structure}. L'ensemble de points $\mu_4 \cup
\{\infty\}$ est stable par $z \mapsto \bar z$, mais cette symétrie renverse
l'orientation du carré, qui fait partie de la structure $\mathrm{I}_4$.}. Les
primes désignent des variantes obtenues par un twist ($\mathrm{II}'_2 =
\mathrm{Tw}_{1/4}\,\mathrm{II}_2$). Le type $\mathrm{I}_6$ est un triangle $J$
et les deux points restants identifiés aux deux orientations de $J$ ; le type
$\mathrm{I}_4$ un point $P$ et un carré orienté sur les quatre autres.

Entre ces sommets, quatorze types de flèches : des $f$ (entre structures de
niveau $2$), des $\gamma$ (« gommage », qui effacent un cercle et montent d'un
niveau) et des $\tau$ (« twist », qui tournent d'une fraction de tour). La page
32 en fait le compte, chaque type étant compté deux fois, par ses sources et par
ses buts :
\[
\begin{array}{lll}
f_{B,s},\ f_{\vec{Q},\alpha} & 60 + 120 & = 180 \\
\tau_{\pi,s},\ \tau^{0}_{r},\ \gamma_{\pi,s},\ \gamma'_{\pi,s},\ \gamma^{0}_{r} &
60 + 120 + 60 + 60 + 120 & = 420 \\
\tau_{\pi,\rho},\ \gamma_{\pi,\rho},\ \tau'_{r},\ \gamma^{1}_{r},\
\gamma^{1\prime}_{r},\ \gamma^{2}_{r},\ \gamma^{2\prime}_{r} & 120 + 120 + 120 +
120 + 360 + 120 + 60 & = 1020
\end{array}
\]
et en tout $1620 = 5 \times 60 + 11 \times 120$ flèches\footnote{La ligne de
$\gamma^{1\prime}_r$ se lit « $3 \times 120 = 3 \times 40 = 3 \times 40$ »,
au-dessus d'une égalité noircie ; c'est $3 \times 120 = 360$ que le total $1020$
exige ($1020 - 660 = 360$), et que donnent ses degrés $3 \to 3$ entre deux types
de $120$ structures. Pour quelques types ($\gamma^{0}_r$, $\gamma^{2}_r$,
$\gamma^{2\prime}_r$), les sources et buts lus sur le tracé du tableau ne
s'accordent pas avec les facteurs du décompte ; on reproduit le décompte. « 1620 »
est écrit sous un « 1580 » biffé.}. Les $542$ sommets forment neuf orbites sous
$\mathfrak{S}_K$, les flèches quatorze.

Les sommets de niveau $0$ se répartissent au-dessus des $15$ décompositions
en pantalons de la sphère à cinq trous, ceux de niveau $1$ au-dessus de ses $10$
cercles ; les décompositions en pantalons sont les sommets du complexe de
Hatcher et Thurston de cette surface, dont les relations élémentaires
comprennent le pentagone\footnote{Rapprochement de cette édition. Hatcher et Thurston ont
publié en 1980 une présentation des groupes modulaires par un complexe de
décompositions en pantalons ; rien dans le dossier ne la cite.}.

Le \emph{programme de travail} de la page 32 :
\begin{enumerate}
\item déterminer les relations fondamentales entre les $1620$ flèches, en les
répartissant par orbites sous $\mathfrak{S}_K$ — en particulier, écrire toutes
les flèches comme composées de flèches « provenant de la dimension modulaire
$\leq 1$ » ;
\item reprendre la question en éliminant les sommets de niveau $0$\footnote{La
page ajoute « (sauf peut-être ceux de type $\mathrm{I}_{10}$ ?) », alors que
$\mathrm{I}_{10}$ est au niveau $2$ ; la suite du passage est encadrée et biffée.} ;
\item éliminer de plus les pièces « tétraédrales » $\mathrm{II}_1$, et toutes
celles qui ne simplifient pas la présentation.
\end{enumerate}
La note encadrée du premier point est le principe des deux étages de
l'\emph{Esquisse} sous forme de tâche à faire.

La page 31 ne porte que des croquis.

\pagerange{33}{35}
\subsection*{Systèmes de cinq points à symétrie (pages 33 à 35, ses pages 1 à 3)}

Soit $\Sigma$ une droite projective sur un corps algébriquement clos $k$ de
caractéristique $0$ (la sphère de Riemann pour $k = \mathbb{C}$), et $K \subset
\Sigma$ de cardinal $5$. On classe les systèmes $(\Sigma, K)$ dont le groupe $G$
des automorphismes n'est pas trivial.

Soit $g \in G$ d'ordre maximal $\nu$. Un automorphisme d'ordre fini $\neq 1$ de
$\Sigma$ a exactement deux points fixes, et $\langle g \rangle$ opère librement
sur leur complémentaire ; donc $K = K^{g} \amalg$ (orbites de cardinal $\nu$),
avec $|K^{g}| \leq 2$, et $2 \leq \nu \leq 5$. Les cas :

\begin{enumerate}
\item[(1)] $\nu = 5$ : $K$ est une orbite, qu'on peut supposer $\mu_5$ ; $G =
\mathbb{D}_5$, et il y a un seul cas à isomorphisme près — la structure
« pentagonale ».
\item[(2)] $\nu = 4$ : $K^{g} = \{P\}$ et $K \setminus \{P\}$ est une orbite ;
on peut prendre $K = \mu_4 \cup \{\infty\}$, $G = \mathbb{Z}/4$, et la catégorie
de ces configurations est équivalente à celle des carrés orientés — la
structure « pyramidale ».
\item[(3)] $\nu = 3$ : $K^{g} = \Sigma^{g}$ a deux points et $J = K \setminus
K^{g}$ est une orbite, qu'on peut prendre égale à $\mu_3$ ; $G = \mathbb{D}_3$,
et la catégorie est équivalente à celle des ensembles $J$ à trois éléments, les
deux pôles correspondant aux deux orientations de $J$ (pour $k = \mathbb{C}$) —
la structure « bitétraédrale ».
\item[(4)] $\nu = 2$ : $K^{g} = \{P\}$ et $I = K \setminus \{P\}$ est réunion de
deux orbites. Alors $g$ est l'une des involutions qui existent pour \emph{tout}
système de quatre points — un élément du groupe de Klein $V(I)$, réalisé par
une homographie — et $P$ un de ses deux points fixes. Cette fois il y a des
modules.
\end{enumerate}

Dans le cas (4), si $I$ est une structure carrée, $P$ est soit un pôle (on
retrouve le cas (2), si $g$ fixe les deux diagonales), soit le milieu de l'une
de deux arêtes opposées du carré sphérique ; et la page demande : « Est-ce le cas
$10$ ? » — la réponse est non : $\mu_4 \cup \{e^{i\pi/4}\}$ n'a qu'un groupe
d'automorphismes d'ordre $2$, et non $\mathbb{D}_5$\footnote{Vérifié
numériquement, en testant les homographies qui envoient trois points de $K$ sur
trois points de $K$. La lecture de « cas $10$ » comme le type $\mathrm{I}_{10}$
de la page 30 est la nôtre.}. Si $I$ est tétraédrale, $P$ est l'un des six
sommets de l'octaèdre associé (les points fixes des trois involutions de $V(I)$).
Le paramétrage de ce cas, barré à la page 34, est repris aux pages 36 à 38.

\emph{Anti-automorphismes} ($k = \mathbb{C}$). Soit $\sigma$ un automorphisme
antiholomorphe de $\Sigma$ qui stabilise $K$. Il n'est pas nécessairement
d'ordre $2$ ni nécessairement doté d'un point fixe : $\sigma(z) = \zeta_5/\bar
z$ stabilise $\mu_5$, est d'ordre $10$ et n'a pas de point fixe. Mais une
puissance impaire de $\sigma$ est une réflexion — une structure réelle — qui
stabilise $K$ : $\sigma$ est conjugué à une isométrie renversant l'orientation
de la sphère ronde, une rotation d'angle $\theta$ composée avec la symétrie par
rapport à l'équateur, dont les orbites ont pour cardinal $2$ (les pôles),
l'ordre $m$ de $\theta$ (sur l'équateur) ou un nombre pair ; comme $5$ est
impair, $K$ rencontre l'équateur et $m$ est impair, et $\sigma^{m}$ est la
symétrie elle-même\footnote{La page affirme que $\sigma$, étant d'ordre pair, a
un point fixe dans $K$, et qu'étant d'ordre fini il est d'ordre $2$. La première
affirmation est fausse (l'exemple ci-dessus) ; la seconde vaut pour un
anti-automorphisme d'ordre fini \emph{qui a un point fixe}. La classification
qui suit n'utilise que l'existence d'une structure réelle stabilisant $K$, que
l'argument ci-dessus fournit.}. On classe donc les structures réelles
$\sigma$ stabilisant $K$ selon $|K^{\sigma}| \in \{1, 3, 5\}$, impair puisque
$\sigma$ est une involution :

\begin{enumerate}
\item[a)] $K^{\sigma} = K$ : cinq points réels, les points réels de
$\mathcal{M}_{0,5}$, de dimension réelle $2$ ; contient le cas pentagonal (avec
$\sigma(z) = 1/\bar z$)\footnote{La page écrit $M^{!}_{0,5} \simeq
\mathrm{Sym}^{*}_{2}(M_{0,3})$, second membre incertain ; on ne le reprend pas.} ;
\item[b)] $|K^{\sigma}| = 3$ : $K = \{0, 1, \infty, \lambda, \bar\lambda\}$,
$\lambda$ dans le demi-plan de Poincaré, dimension réelle $2$ ; contient le cas
bitétraédral ;
\item[c)] $K^{\sigma} = \{P\}$ : $I = K \setminus \{P\}$ est un système de
quatre points stable par une conjugaison sans point fixe sur $I$, de dimension
modulaire réelle $1$, et le point réel $P$ ajoute une dimension : dimension
réelle $2$ ; contient le cas pyramidal (avec $\sigma(z) = i\bar z$).
\end{enumerate}

La page 35 range enfin les strates de $\overline{\mathcal{M}}_{0,5}$ par dimension
modulaire, avec le graphe de chaque strate (un sommet à cinq pattes ; deux
sommets ; trois sommets en chaîne) :
\[
\begin{array}{llll}
\text{dimension modulaire} & 2 & 1 & 0 \\
\text{nombre de composantes} & 1 & 10 & 15 \\
\text{réalisations spéciales} & 12 & 10 \cdot 3 \cdot 2 + 10 \cdot 2 \cdot 6 =
60 + 120 & 15 \cdot 2 \cdot 2 = 60
\end{array}
\]
et note, en dimension $2$, $5 \cdot 3 \cdot 2 = 30$ structures pyramidales et
$10 \cdot 2 = 20$ bitétraédrales. Ces nombres sont ceux des types non primés du
tableau de la page 30 ($\mathrm{I}_{10}$ ; $\mathrm{II}_2$ et $\mathrm{II}_1$ ;
$\mathrm{III}_2$) ; les types primés, obtenus par twist, s'y
ajoutent\footnote{Ce rapprochement des deux décomptes est le nôtre ; le dossier
ne les confronte pas.}.

\pagerange{36}{38}
\subsection*{Le cas involutif, paramétré (pages 37, 36, 38)}

Les trois feuillets se lisent dans l'ordre 37, 36, 38\footnote{L'ordre des
archivistes, 36, 37, 38, est gardé pour les pages ; c'est le contenu qui
impose la lecture : la page 37 s'arrête au milieu d'une phrase que la page 36
continue, et la page 38 suit la page 36.}. On reprend le cas (4) : $K = \{P\}
\amalg I$, avec une involution holomorphe $u$ de $\Sigma$ stabilisant $I$, sans
point fixe dans $I$, et fixant $P$.

\emph{La donnée.} Une sphère $\Sigma$ munie d'une involution $u$ et d'un point
fixe marqué, appelé $\infty$ (l'autre point fixe étant $0$), équivaut à un
torseur $T_1 = \Sigma \setminus \{0, \infty\}$ sous $\mathbb{C}^{*}$, c'est-à-dire
à une droite vectorielle complexe $L_1$ avec $T_1 = L_1 \setminus \{0\}$ ; $u$ est
$x \mapsto -x$. Une paire d'orbites de $\{\pm 1\}$ dans $T_1$ équivaut, via $x
\mapsto x^{\otimes 2}$, à une paire $\{\xi, \eta\}$ de points distincts de $T_2 =
L_2 \setminus \{0\}$, $L_2 = L_1^{\otimes 2}$. Le produit $e = \xi \eta \in T_4$
($L_4 = L_1^{\otimes 4}$) trivialise $T_4$, d'où une réduction du groupe structural
de $T_1$ de $\mathbb{C}^{*}$ à $\mu_4$ : les bases $e_1$ de $L_1$ telles que
$e_1^{\otimes 4} = e$.

\emph{Les invariants.} Une telle base choisie, $\xi = \xi_0 e_1^{2}$ et $\eta =
\eta_0 e_1^{2}$ avec $\xi_0 \eta_0 = 1$ ; la donnée équivaut à
\[
\mu = \xi_0 + \eta_0 \in \mathbb{C}, \qquad \mu \neq \pm 2,
\]
$\xi_0$, $\eta_0$ étant les racines de $z^2 - \mu z + 1 = 0$, distinctes
exactement quand $\mu^2 \neq 4$. On retrouve $I$ comme l'ensemble des racines de
$z^4 - \mu z^2 + 1 = 0$,
\[
I = \Bigl\{ \pm\sqrt{\tfrac{1}{2}\bigl(\mu \pm \sqrt{\mu^2 - 4}\bigr)} \Bigr\}
= \Bigl\{ \pm\sqrt{\mu' \pm \sqrt{\mu'^2 - 1}} \Bigr\}, \qquad \mu' =
\tfrac{1}{2}\mu .
\]
Sans choisir de base, l'invariant est
\[
\lambda = \frac{(\xi + \eta)^2}{\xi\eta} = \frac{\xi}{\eta} + \frac{\eta}{\xi} +
2 = \mu^2, \qquad \lambda \neq 4,
\]
c'est-à-dire $\lambda - 2 = \xi/\eta + \eta/\xi \neq 2$\footnote{La page écrit
$\xi \wedge \eta$ pour le produit $\xi\eta \in T_4$ (produit contracté des
torseurs), et ne relève pas que $\lambda = \mu^2$ ; l'égalité est immédiate
puisque $\xi_0\eta_0 = 1$.}.

\emph{Le changement de base.} Remplacer $e_1$ par $\zeta e_1$ ($\zeta^4 = 1$)
change $(\xi_0, \eta_0)$ en $(\zeta^2\xi_0, \zeta^2\eta_0)$, donc $\mu$ en
$\zeta^2\mu = \pm\mu$. L'espace des modules des structures rigidifiées par le
choix de $e_1$ est donc un ouvert de la droite des $\mu$, qui s'envoie sur
$\mathcal{M}_{0,5}$ en passant au quotient par $\mu \mapsto -\mu$ ; la page
l'écrit $\mathbb{P}^{1} \setminus \{0, \pm 2, \infty\}$ et propose de regarder
l'effet sur les groupoïdes fondamentaux\footnote{L'exclusion de $\mu = 0$ n'est
pas justifiée sur la page. Pour $\mu = 0$, $I$ est l'ensemble des racines de
$z^4 = -1$, et $K = I \cup \{\infty\}$ est la configuration pyramidale du cas
(2), où le groupe d'automorphismes devient $\mathbb{Z}/4$ ; c'est aussi le point
fixe de $\mu \mapsto -\mu$. La raison est vraisemblablement celle-là ; elle est
de cette édition.}.

\pagerange{39}{42}
\section*{Le catalogue des opérations et la tour (pages 39 à 42)}

La page 39, une couverture, porte « Catalogue des opérations élémentaires dans
$R\mathcal{T}D_{\mathrm{top}}$ »\footnote{La première lettre est lue $R$ sans
certitude ; la page 42 écrit le même sigle. Rien ne dit si ce $R$ est celui de
$SRT(X)$, page 1.}.

\pagerange{40}{40}
\subsection*{La première liste (page 40, octobre 1983)}

La page 40 fixe la donnée combinatoire d'une surface découpée et rigidifiée :
un ensemble $\vec{A}$ de demi-arêtes avec une involution $\sigma$, un ensemble
$RR$ de « repères » sur lequel opère le groupe diédral infini
$\mathbb{D}_\infty = \langle \sigma_0, \sigma_1 \mid \sigma_0^2 = \sigma_1^2 =
1 \rangle$, une application $\mathrm{supp} : RR \to \vec{A}$, un ensemble $S$ de
sommets et un genre $g : S \to \mathbb{N}$, avec les conditions :
\begin{enumerate}
\item[a)] l'involution est d'ordre $2$ et « sans point fixe »\footnote{La page
écrit « $\vec{\sigma}^{2} = \mathrm{id}$, $\sigma$ sans pt fixe », « sans »
incertain. Les arêtes libres étant des points fixes de $\sigma$ sur $\vec{A}$,
la condition ne peut porter que sur les arêtes liées ou sur l'action sur $RR$ ;
on ne tranche pas.} ;
\item[b)] $\mathrm{supp}$ est compatible avec $\mathbb{D}_\infty \to \{\pm 1\}$,
$\{\pm 1\}$ opérant sur $\vec{A}$ par $\sigma$ ;
\item[c)] $RR/\mathbb{D}_\infty \xrightarrow{\sim} A = \vec{A}/\sigma$ ;
\item[d)] éventuellement, la condition anabélienne : pour tout sommet $s$, $2g_s
+ \nu_s \geq 3$, où $\nu_s$ est le nombre de demi-arêtes en $s$.
\end{enumerate}
C'est le « graphe pondéré à repères » des pages 49 à 51, sous une forme déjà
complète\footnote{La page écrit aussi $3g_s - 3 + \nu_s \geq 0$ et $2g_s - 2 +
\nu_s \geq 1$. La seconde équivaut à $2g_s + \nu_s \geq 3$ ; la première est
plus faible : elle admet le tore sans trou, $(g, \nu) = (1, 0)$, que les deux
autres excluent.}.

Puis la liste des opérations : somme ; amalgamation (sans gommage) et découpage
(en gardant les données de recollement), « inverses l'une de l'autre » ; gommage
de sommets de marquage et surmarquage, groupés de même ; bouchage des trous
(gommage des cercles de raccord) ; gommage des cercles de découpage. Le
vocabulaire se fixe : cercles de découpage (intérieurs) et cercles de raccord
(le bord), sommets de marquage, ou de rigidification, cercles marqués ou non ; et
les types combinatoires, rangés en un losange
\[
\begin{array}{ccc}
 & \mathcal{T} & \\
S\mathcal{T} & & D\mathcal{T} \\
 & SD\mathcal{T} &
\end{array}
\]
— sans structure, avec marquage $S$, avec découpage $D$, avec les deux —,
annotés « graphes pondérés discrets », « graphes b. », « graphes pondérés »,
« graphes ».

\emph{Le twist relatif.} Pour un cercle de découpage marqué et une fraction de
tour $\rho$, le twist $\tau_\rho$ se définit par découpage puis recollement
après une rotation de $\rho$ tour. Pour $\rho \in \mathbb{Z}$, l'objet recollé
est le même ; mais la flèche $\tau_\rho$ n'est pas l'identité : c'est la
puissance $\rho$-ième du twist de Dehn le long du cercle\footnote{La page écrit
« $\tau_\rho = \mathrm{id}$ si $\rho \in \mathbf{Z}$ », vrai de l'opération sur
les objets, faux des flèches du groupoïde. Ce sont les $\mathbb{R}^{I}$ et
$\mathbb{Z}^{I}$ des pages 1 à 3, en un cercle intérieur.}. Enfin, encadrée,
la ligne
\[
S_{*}\mathcal{T}^{\pm}_{g,\nu} \longrightarrow \mathcal{T}^{\pm}_{g,\nu}
\longrightarrow \mathrm{Ens}(\nu) \times [\pm 1]
\]
qu'on retrouvera page 61.

\pagerange{41}{42}
\subsection*{La seconde liste, et les douze avatars (pages 41 et 42)}

La page 41 reprend le catalogue « opérations $+$ structures » : somme, gommage
(partiel) du marquage et marquage ; découpage (partiel) et amalgamation ou
recollage, appariés ; bouchage de trous (gommage des cercles de raccord) ;
gommage des cercles internes ; avec une condition « anabélienne » ou « de
stabilité », des foncteurs et des isomorphismes de « twist ». Deux losanges, le
second marqué « types combinatoires », sont reliés par des flèches du premier
vers le second\footnote{La lettre du second losange est lue $\mathbb{L}$ sans
certitude.}.

La page 42 énumère « les douze avatars principaux de la Tour de Teichmüller »,
une colonne de groupoïdes où chacun s'envoie sur le précédent :
\[
\mathcal{T}_{\mathrm{top}} \leftarrow \mathcal{T}_{\mathrm{diff}} \leftarrow
\mathcal{T}_{\mathrm{an}} \leftarrow \mathcal{T}_{\mathrm{hyp}} \leftarrow
\mathcal{T}_{\mathbb{C}\text{-alg}} \leftarrow
\mathcal{T}_{\overline{\mathbb{Q}}\text{-alg}} \leftarrow
\mathcal{T}_{\mathrm{comb}} \leftarrow \mathcal{T}_{\mathrm{Lego}} \leftarrow
\mathcal{T}_{\mathbb{P}}
\]
— surfaces topologiques, différentiables (à bord lisse), $\mathbb{C}$-analytiques
(à bord analytique), hyperboliques (à bord géodésique), courbes algébriques
complexes de Mumford–Deligne (avec rigidification aux points doubles et points
marqués), courbes définies sur $\overline{\mathbb{Q}}$, cartes orientées à bord
avec sommets marqués, jeu de « Légo » à pièces combinatoires, jeu
« pantalonnique » dont les seules pièces sont les triades —, plus trois variantes
$\mathcal{T}'_{\mathrm{diff}}$, $\mathcal{T}'_{\mathrm{an}}$,
$\mathcal{T}'_{\mathrm{hyp}}$ à points marqués et vecteurs tangents, chacune
équivalente à la non primée. Le compte fait douze\footnote{La page dessine le
$T$ à double jambage ; on l'écrit $\mathcal{T}$. Le nombre de pièces du jeu de
Légo est lu « 10 » sans certitude. Le passage de $\mathcal{T}_{\mathbb{C}\text{-alg}}$
à $\mathcal{T}_{\overline{\mathbb{Q}}\text{-alg}}$ est celui où le groupe de Galois
de $\overline{\mathbb{Q}}$ agit sur la tour ; c'est le lien de ce dossier avec la
\emph{Longue Marche}.}.

\pagerange{43}{60}
\section*{Groupoïdes de Teichmüller (pages 43 à 60, ses pages 1 à 17)}

\pagerange{43}{45}
\subsection*{Les groupoïdes topologiques}

Le \emph{groupoïde de Teichmüller orienté topologique}
$\mathcal{T}^{+}_{\mathrm{top}}$ a pour objets les surfaces compactes orientées
à bord, pour flèches les classes d'isotopie d'homéomorphismes conservant
l'orientation. Le \emph{biorienté} $\mathcal{T}^{\pm}_{\mathrm{top}}$ a les
mêmes objets et toutes les classes d'isotopie d'homéomorphismes. Sur les
surfaces connexes, la signature est un morphisme $\chi :
\mathcal{T}^{\pm}_{\mathrm{top}} \to [\pm 1]$ dont la fibre est
$\mathcal{T}^{+}_{\mathrm{top}}$ ; sur une surface non connexe un homéomorphisme
peut conserver l'orientation sur une composante et la renverser sur une autre,
et il n'y a plus de signature — ce que dit la marge\footnote{« canulé, si on ne
se borne à des surfaces connexes ! », lecture incertaine. Une autre note, dans
l'angle de la page, propose d'appeler $\mathcal{T}$ le biorienté, « le plus
utilisé », et $\mathcal{T}^{+}$ l'orienté ; d'où la convention de cette lecture,
qui écrit toujours l'exposant.}.

La variante la plus complète est le \emph{groupoïde de Teichmüller spécial
divisé} $S_{*}\mathcal{T}D^{+}_{\mathrm{top}}$ : ses objets sont les triples
$(X, \widetilde{\Sigma}, R)$ où $X$ est une surface compacte orientée à bord,
$\widetilde{\Sigma} = \partial X \sqcup \Sigma$ une sous-variété compacte de
dimension $1$ contenant $\partial X$, et $R \subset \widetilde{\Sigma}$ une
partie finie ; ses flèches, les classes d'isotopie d'homéomorphismes conservant
l'orientation et respectant $\widetilde{\Sigma}$ et $R$. Les composantes de
$\partial X$ sont les \emph{cercles de raccord} (ou de recollement, ou « trous
»), celles de $\Sigma$ les \emph{cercles de découpage}, les points de $R$ les
\emph{sommets de rigidification}\footnote{La phrase de la page 45 est remaniée en
interligne et l'ordre de ses morceaux est incertain ; l'attribution des deux
noms se lit sans ambiguïté aux pages 46, 50 et 57.}. Le biorienté
$S_{*}\mathcal{T}D^{\pm}_{\mathrm{top}}$ se définit de même, avec $S_{*}\mathcal{T}
D^{+}_{\mathrm{top}} \hookrightarrow S_{*}\mathcal{T}D^{\pm}_{\mathrm{top}} \to
[\pm 1]$. Il contient les sous-groupoïdes pleins
\begin{itemize}
\item $R = \Sigma = \emptyset$ : $\mathcal{T}^{+}_{\mathrm{top}}$,
$\mathcal{T}^{\pm}_{\mathrm{top}}$ ;
\item $\Sigma = \emptyset$ : $S_{*}\mathcal{T}^{+}_{\mathrm{top}}$,
$S_{*}\mathcal{T}^{\pm}_{\mathrm{top}}$, les groupoïdes \emph{spéciaux} ou
rigidifiés ;
\item $R = \emptyset$ : $\mathcal{T}D^{+}_{\mathrm{top}}$,
$\mathcal{T}D^{\pm}_{\mathrm{top}}$, « pas intéressant je crois, sauf pour
$\Sigma = \emptyset$ également ».
\end{itemize}
Pour une surface connexe de type $(g, \nu)$, les groupes d'automorphismes sont
les groupes modulaires d'aujourd'hui : dans $\mathcal{T}^{+}_{\mathrm{top}}$, le
groupe modulaire qui permute les bords ; dans $S_{*}\mathcal{T}^{+}_{\mathrm{top}}$
avec un point sur chaque bord, le groupe $S_{*}T(X, \underline{s})$ des pages 1
à 3. Les deux lignes barrées de la page 44 sont un premier départ de la phrase
de la page 51.

\pagerange{45}{48}
\subsection*{Le graphe de Mumford--Deligne}

Un \emph{graphe bipondéré} est une donnée $(S, \vec{A}, \sigma, o, \gamma, r)$ :
$S$ et $\vec{A}$ finis, $\sigma$ une involution de $\vec{A}$ dont les points
fixes $\vec{A}_\ell = \vec{A}^{\sigma}$ sont les \emph{arêtes libres}, $o :
\vec{A} \to S$, un genre $\gamma : S \to \mathbb{N}$ et un poids $r : A \to
\mathbb{N}$ sur les arêtes $A = \vec{A}/\sigma = A_0 \sqcup A_\ell$. Sommets
isolés, boucles, arêtes multiples sont permis. De façon équivalente, c'est un
graphe ordinaire $(S, \vec{A}_0, \sigma, o)$, plus un ensemble fini $A_\ell$
et une application $A_\ell \to S$\footnote{La page commence par donner la
pondération aux seuls sommets ; deux notes de la marge disent que c'est « pas
assez précis » et que la pondération d'une arête devrait être celle d'un polygone
combinatoire $T_a$ avec $\omega(T_a) \simeq \omega(a)$ — ce que font les pages 49
et 50. Un bloc barré en croix posait $\widetilde{S} = S \sqcup A_\ell$.}. Avec
leurs isomorphismes, ils forment un groupoïde (Gr.bip).

Le foncteur
\[
(*)\qquad \Gamma : S_{*}\mathcal{T}D^{\pm}_{\mathrm{top}} \longrightarrow
(\mathrm{Gr.bip})
\]
associe à $(X, \widetilde{\Sigma}, R)$ le graphe dont
\begin{itemize}
\item les sommets sont les composantes de $X \setminus \widetilde{\Sigma}$ (les
« faces ») ;
\item les demi-arêtes sont les couples $(C, \text{rive})$, $C \in
\pi_0(\widetilde{\Sigma})$ et une rive de $C$ — deux pour un cercle de découpage,
une seule (celle qui est dans $X$) pour un cercle du bord ;
\item $\sigma$ échange les deux rives d'un cercle de découpage, et fixe la rive
unique d'un cercle du bord ;
\item $o$ envoie une rive sur la face qui la contient ;
\item $\gamma$ est le genre de (l'adhérence de) chaque face ;
\item $r(C) = \mathrm{card}(R \cap C)$\footnote{La page définit une demi-arête
comme un couple $(C, \omega)$, $\omega$ une orientation de $C$, identifiée à une
rive par $\mathrm{Or}(C) \simeq \mathrm{rive}(C)$. Cette identification dépend de
l'orientation de $X$, que les flèches de $\mathcal{T}^{\pm}$ peuvent renverser ; on
prend donc les rives. Le poids, noté $\rho$ dans le sextuple, devient $r$ à la
dernière ligne.}.
\end{itemize}
C'est le graphe dual d'une courbe stable, ce qu'on appelle aujourd'hui un
\emph{graphe stable} ou graphe modulaire, avec ses pattes (les arêtes libres) ;
la marge le nomme : « MD = modulaire Mumford–Deligne ».

\emph{Proposition (page 47).} Le foncteur $(*)$ induit une bijection sur les
$\pi_0$ et une surjection sur les groupes d'automorphismes\footnote{La page dit
« c'est un $0$-isomorphisme, en langage homotopique », le $0$ incertain. Dans le
vocabulaire d'aujourd'hui, un foncteur entre groupoïdes bijectif sur les $\pi_0$
et surjectif sur les $\pi_1$ est \emph{$1$-connexe}. « est surjectif et » est
biffé.}. Bijection : une surface compacte orientée est classée par son genre et
son nombre de bords, l'objet se reconstruit en recollant des surfaces de type
$(\gamma(s), \text{valence de } s)$ selon le graphe, et deux parties finies de
même cardinal d'un cercle sont isotopes. Surjection : tout automorphisme du
graphe pondéré se réalise par un homéomorphisme. Le même argument montre que
$\pi_0(S_{*}\mathcal{T}D^{+}_{\mathrm{top}}) \to
\pi_0(S_{*}\mathcal{T}D^{\pm}_{\mathrm{top}})$ est bijectif, le graphe ne voyant pas
l'orientation.

\emph{Remarques (pages 47 et 48).} Les deux groupoïdes ont une opération
évidente de \emph{somme} (réunion disjointe), et $(*)$ la respecte ; c'est la
plus triviale des opérations à transcrire combinatoirement. On a $\pi_0(X)
\simeq \pi_0(\Gamma(X, \widetilde{\Sigma}, R))$, et le genre d'une composante
connexe se lit sur son graphe :
\[
g = \sum_{s \in S} \gamma(s) + \operatorname{rg} H^{1}(\Gamma),
\]
le second terme comptant les cycles indépendants du graphe (les pattes n'y
contribuent pas).

\emph{Cas particuliers.} Pour $S_{*}\mathcal{T}^{\pm}_{\mathrm{top}}$ ($\Sigma =
\emptyset$), les graphes n'ont que des arêtes libres : des sommets isolés munis
de pattes, donnés par $\gamma : S \to \mathbb{N}$, $A_\ell \to S$ et $r : A_\ell
\to \mathbb{N}$. Pour $\mathcal{T}^{\pm}_{\mathrm{top}}$, de plus $r = 0$. Dans
le cas connexe, le graphe est un sommet de genre $g$ portant $\nu$ pattes, et
l'objet est classé par le couple $(g, \nu)$.

\pagerange{49}{51}
\subsection*{Graphes à repères}

Les remarques de la page 49 partent de là : le poids $r(a)$ ne retient que le
nombre de repères sur le cercle $C_a$, et la structure est plus riche. Notons $\omega(a)$ l'ensemble des demi-arêtes
au-dessus de $a$ (une ou deux) et $R_a = R \cap C_a$. Chaque $\omega \in
\omega(a)$ — une rive, donc une orientation de $C_a$ — munit $R_a$ d'un ordre
circulaire $u_\omega$, et $u_{\omega'} = u_\omega^{-1}$ si $\omega' \neq \omega$.
De façon équivalente, si $r(a) \geq 1$, $R_a$ est un torseur sous
\[
\mathbb{Z}_a/r(a), \qquad \mathbb{Z}_a = \mathbb{Z}
\wedge^{\{\pm 1\}} \omega(a),
\]
le groupe $\mathbb{Z}$ tordu par l'ensemble $\omega(a)$ des orientations (pour
une arête libre, $\omega(a)$ est un singleton et $\mathbb{Z}_a = \mathbb{Z}$)
\footnote{Pour $r(a) = 0$, $R_a$ est vide, et n'est pas un torseur ; la page ne
le distingue pas ici, mais les conditions $r_a \neq 0$ ou $r_a \geq 1$ de la page
60 en tiennent compte.}.

Pour chaque $a$, on dispose ainsi d'un polygone combinatoire $P_a$ de sommets
$R_a$ — même si $r(a) = 1$ ou $2$ — et la famille des $(R_a)$ peut être
remplacée par un ensemble $\vec{R}$ de « repères de rigidification », muni d'une
action de $\mathbb{D}_\infty = \langle \sigma_0, \sigma_1 \mid \sigma_0^2 =
\sigma_1^2 = 1\rangle$ et d'une application $\vec{R} \to \vec{A}$ compatible avec
$\varepsilon : \mathbb{D}_\infty \to \{\pm 1\}$, plus une partie $\vec{R}^{+}_\ell$
des repères au-dessus des arêtes libres qui oriente les cercles du bord. On peut
prendre pour $\vec{R}$ l'ensemble des \emph{drapeaux} des polygones $P_a$ —
couples (sommet, côté incident) —, $\sigma_0$ changeant le sommet et $\sigma_1$
le côté ; un drapeau définit un sens de parcours de $C_a$, donc une rive, et
chaque $\sigma_i$ le renverse, d'où $\varepsilon(\sigma_0) = \varepsilon(\sigma_1)
= -1$\footnote{L'interprétation de $\vec{R}$ par les drapeaux, et les valeurs de
$\varepsilon$, sont les nôtres ; la page ne les donne pas. Elles sont celles qui
rendent vraie la condition c) de la page 40, $RR/\mathbb{D}_\infty \simeq A$, une
orbite par arête marquée. $\mathbb{D}_\infty$ est le groupe cartographique en
dimension $1$, comme le groupe à trois générateurs $\sigma_0, \sigma_1, \sigma_2$
du dossier 141 (pages 9 et 10) l'est en dimension $2$.}.

On obtient la notion de \emph{graphe pondéré à repères},
\[
\vec{R}^{+}_\ell \hookrightarrow \vec{R} \longrightarrow \vec{A}
\xrightarrow{\ o\ } S \xrightarrow{\ \gamma\ } \mathbb{N},
\]
et un foncteur $S_{*}\mathcal{T}D^{\pm}_{\mathrm{top}} \to (\mathrm{Gr.pondrep})$
« qui sera important sans doute dans l'explicitation de l'opération de
recollement ». Toute description algébrico-combinatoire de
$S_{*}\mathcal{T}D^{\pm}_{\mathrm{top}}$ devra permettre de reconstituer, à
isomorphisme unique près, ce foncteur — et non seulement le foncteur vers les
graphes bipondérés.

\pagerange{51}{59}
\subsection*{Les opérations essentielles}

\textbf{A) Gommages.}

\emph{(1) Gommage de repères.} À $(X, \Sigma, R)$ muni d'une partie $R' \subset
R$, on associe $(X, \Sigma, R \setminus R')$. Une description combinatoire devra
donc fournir aussi le groupoïde des objets munis d'une partie marquée $R'$, et
un foncteur « gommage de $R'$ », s'inscrivant dans un carré commutatif
\[
\begin{array}{ccc}
S_{*}\mathcal{T}D^{\pm}_{\mathrm{top}}(R' \text{ marqué}) &
\xrightarrow{\ \text{gommage de } R'\ } & S_{*}\mathcal{T}D^{\pm}_{\mathrm{top}} \\
\downarrow{\scriptstyle \text{oubli de } R'} & & \downarrow \\
S_{*}\mathcal{T}D^{\pm}_{\mathrm{top}} & \longrightarrow & (\mathrm{Gr.bip})/r
\end{array}
\]
où le coin est le groupoïde des graphes pondérés par $\gamma$ seulement. Au
niveau des graphes à repères, l'opération enlève une partie $\vec{R}'$ de
$\vec{R}$ stable par $\mathbb{D}_\infty$\footnote{Une longue note oblique de la
marge de la page 52, presque illisible, remarque que l'opération, qui semble un
affaiblissement de structure, est plutôt, à certains égards, un
\emph{renforcement} ; la page 56 le dit pour le gommage partiel.}.

\emph{(2) Gommage de cercles de découpage.} On se donne une réunion $\Sigma'$ de
composantes de $\Sigma$, c'est-à-dire une partie $A'_0$ des arêtes liées, et
l'on remplace $\Sigma$ par $\Sigma \setminus \Sigma'$ ; les repères situés sur
$\Sigma'$ sont gommés avec. Sur les graphes, c'est la \emph{contraction} des
arêtes de $A'_0$, et le graphe ancien s'envoie sur le nouveau par un
« morphisme de généralisation (à l'opposé de spécialisation) » : on remonte
d'une strate du bord de $\overline{\mathcal{M}}_{g,\nu}$ vers une strate moins
dégénérée\footnote{La page écrit $\pi_0(\Sigma')$ là où l'on attend
$\pi_0(\Sigma)$.}. La pondération « suit », et la formule du genre de la page
48 dit comment : chaque sommet du nouveau graphe, image d'un sous-graphe connexe
$\Gamma'$ contracté, a pour genre $\sum_{s \in \Gamma'} \gamma(s) +
\operatorname{rg} H^{1}(\Gamma')$\footnote{La page s'arrête sur « Il convient ici
d'expliciter comment la pondération suit\ldots » ; la formule est la nôtre,
déduite de celle de la marge de la page 48.}.

\emph{(3) Bouchage de trous.} Pour une partie $A'_\ell$ des cercles du bord, on
bouche les trous correspondants par des disques (des cônes sur ces cercles). Sur
les graphes : effacer les pattes de $A'_\ell$ et les repères qui sont au-dessus ;
les genres des sommets ne changent pas.

\emph{Gommage simultané mixte.} Pour une partie ouverte et fermée $R' \sqcup
\Sigma'$ de $R \sqcup \Sigma$ telle que $R' \supset R \cap \Sigma'$, on fait les
deux à la fois : enlever $\vec{R}'$, contracter les arêtes de $\Sigma'$, faire
suivre la pondération\footnote{La page écrit « contracter en un point les arêtes
de $A'_\ell$ (considérées comme libres que liées) », l'indice lu sans certitude ;
il s'agit des arêtes de $\Sigma'$.}.

\textbf{B) Marquage.} Soit $A^{+}$ l'ensemble des arêtes effectivement marquées
($r(a) > 0$), $\Sigma^{+}$ la réunion des cercles correspondants et $B^{+}$
l'ensemble des composantes de $\Sigma^{+} \setminus R$, des intervalles ouverts.
Un \emph{surmarquage} est donné par $r' : B^{+} \to \mathbb{N}$ : on ajoute
$r'(b)$ nouveaux repères dans chaque intervalle $b$. Il y a des choix, mais
l'objet obtenu $(X, \Sigma, R_1)$, muni de la partie $R_1 \setminus R$, est
déterminé à isomorphisme unique près — l'espace des positions de $r'(b)$ points
dans un intervalle est contractile — et le gommage de $R_1 \setminus R$ rend
l'objet de départ.

On se bornera sans doute à des gommages et surmarquages \emph{uniformes}, qui se
lisent sur les torseurs $R_a$ sous $\mathbb{Z}_a/r(a)$ :
\begin{itemize}
\item le \emph{gommage partiel uniforme} garde, pour un diviseur $r'(a)$ de
$r(a)$, l'une des $d(a) = r(a)/r'(a)$ orbites du sous-groupe d'ordre $r'(a)$ de
$\mathbb{Z}_a/r(a)$ : c'est une restriction du groupe structural à ce
sous-groupe, isomorphe à $\mathbb{Z}_a/r'(a)$ ;
\item le \emph{surmarquage uniforme} d'exposant $d(a)$ remplace $R_a$ par
$R_a \wedge^{\mathbb{Z}_a/r(a)} \mathbb{Z}_a/d(a)r(a)$, le long de l'inclusion
$\mathbb{Z}_a/r(a) \hookrightarrow \mathbb{Z}_a/d(a)r(a)$ : on intercale
$d(a) - 1$ points entre deux repères consécutifs.
\end{itemize}
Le premier demande un choix (l'orbite gardée) : c'est un \emph{renforcement} de
structure ; le second n'en demande aucun et oublie quels points étaient les
anciens : c'est un affaiblissement\footnote{La page écrit « restriction des
groupes d'opérateurs de $\mathbf{Z}_a/r(a)$ à $\mathbf{Z}_a/r'(a)$ » et
« extension \ldots{} à $\mathbf{Z}_a/r''(a)$ » ; $\mathbb{Z}/r'$ n'étant pas un
sous-groupe mais un quotient de $\mathbb{Z}/r$, on lit : restriction au
sous-groupe d'ordre $r'$, extension le long de l'inclusion. Ces opérations sont
les twists fractionnaires des pages 1 à 3 vus du côté des repères.}.

\textbf{C) Découpage}, le long d'une réunion $\Sigma'$ de cercles de découpage,
c'est-à-dire d'une partie $A'$ des arêtes liées : sur les graphes, couper chaque
arête par son milieu en deux pattes. Le marquage suit. Si le cercle est
effectivement marqué ($C \cap R \neq \emptyset$), l'opération s'inverse à une
ambiguïté finie près — le choix de l'identification des deux torseurs de
repères ; sinon, elle ne s'inverse pas. Dans la façon dont on va travailler,
dit-il, les cercles de découpage seront toujours effectivement marqués ; mais on
veut garder le droit de les gommer.

\textbf{D) Recollement}, « la plus essentielle, la moins triviale peut-être de
toutes ces opérations ». On se donne
\begin{enumerate}
\item[a)] une partie $A'$ des arêtes libres et une involution sans point fixe
$\sigma'$ de $A'$ ;
\item[b)] pour chaque paire $\{a, a'\}$, des homéomorphismes renversant
l'orientation $C_a \to C_{a'}$ et $C_{a'} \to C_a$, inverses l'un de l'autre.
\end{enumerate}
Pour que l'opération soit définie à isomorphisme \emph{unique} près dans
$S_{*}\mathcal{T}D^{\pm}_{\mathrm{top}}$, il faut que les cercles recollés soient
effectivement marqués, $A' \subset A^{+}$ ; quitte à gommer ou surmarquer
uniformément au préalable, on suppose $r(a) = r(a')$, et l'on se donne un
anti-isomorphisme des torseurs $R_a$ et $R_{a'}$ sous $\mathbb{Z}/r$,
c'est-à-dire un élément de $R_a \wedge_{\mathbb{Z}/r} R_{a'}$. Le recollement
est alors bien déterminé : les homéomorphismes renversant l'orientation qui
réalisent un anti-isomorphisme donné $R_a \to R_{a'}$ forment un espace
contractile, alors que sans repères ils forment un espace qui a le type
d'homotopie d'un cercle, et le résultat n'est déterminé qu'à twist de Dehn
près\footnote{Cette justification est la nôtre ; la page énonce les conditions
sans dire pourquoi elles suffisent.}. Sur les graphes : rattacher deux à deux
les pattes par leurs bouts ; les genres des sommets ne changent pas, poids et
repères suivent.

Le recollement et la somme sont, en termes d'aujourd'hui, les compositions d'une
\emph{opérade modulaire} : le recollement de deux bords d'une même composante ou
de deux composantes distinctes, indexé par des graphes stables\footnote{Le mot
est de Getzler et Kapranov (1998) ; rapprochement de cette édition.}.

\pagerange{60}{60}
\subsection*{La récapitulation (page 60, sa page 17)}

\begin{itemize}
\item[$0^\circ$] somme directe ;
\item[$1^\circ$] gommage total des repères d'un cercle, $r_a \mapsto 0$ (on
suppose $r_a \neq 0$) — fortement irréversible ;
\item[$2^\circ$] gommage partiel uniforme d'exposant $d_a = r_a/r'_a > 1$ —
renforcement de structure faible ;
\item[$3^\circ$] surmarquage uniforme d'exposant $d_a > 1$ ($r_a \neq 0$) —
affaiblissement de structure faible ;
\item[$4^\circ$] gommage d'un cercle de découpage $C_a$, $a \in A_0$ — fortement
irréversible ;
\item[$5^\circ$] bouchage d'un trou $C_a$, $a \in A_\ell$ — fortement
irréversible ;
\item[$6^\circ$] découpage le long de $C_a$, $a \in A_0$ — réversible à un choix
fini près si $r_a \geq 1$ ;
\item[$7^\circ$] recollement le long de $C_a$, $C_{a'}$ ($a \neq a'$ libres,
$r_a = r_{a'} \geq 1$) via un élément de $R_a \wedge R_{a'}$ — réversible à un
choix fini près.
\end{itemize}
Pour $2^\circ$, il y a $d_a$ choix possibles ; pour $6^\circ$, l'opération est
fortement irréversible si $r_a = 0$, et la page se demande s'il faut séparer les
deux cas ou exclure $r_a = 0$.

\pagerange{61}{64}
\section*{Le programme de dévissage (pages 61 à 64, ses pages 18 à 21)}

Il faudrait, dit la page 61, un « petit théorème de dévissage » : à équivalence
près, $S_{*}\mathcal{T}D^{\pm}_{\mathrm{top}}$, avec ses foncteurs vers les
graphes pondérés et vers $[\pm 1]$ et les opérations élémentaires $1^\circ$ à
$7^\circ$ — sauf $4^\circ$ et $5^\circ$ —, se reconstitue à partir des
groupoïdes connexes
\[
S_{1}\mathcal{T}^{\pm}_{g,\nu} \longrightarrow \mathcal{T}^{\pm}_{g,\nu}
\longrightarrow \mathrm{Ens}(\nu) \times [\pm 1]
\]
\footnote{L'indice sous le $S$, lu $f$, est douteux ; on y lit l'étoile de
$S_{*}$. La marge ajoute que « ça ne suffit pas tout à fait » pour avoir toutes
les opérations, notamment $4^\circ$ et $5^\circ$.}. Ici
$S_{1}\mathcal{T}^{\pm}_{g,\nu}$ est le sous-groupoïde des $(X, R)$ avec $X$
connexe de genre $g$, à $\nu$ bords, et un repère sur chaque bord — le graphe
est une étoile, un sommet de genre $g$ et $\nu$ pattes de poids $1$ —, et
$\mathcal{T}^{\pm}_{g,\nu}$ celui des surfaces sans structure de type $(g, \nu)$.
La première flèche oublie les repères : sur les groupes, c'est $S_{*}T(X,
\underline{s}) \to T(X)$ des pages 1 à 3, de noyau les twists de bord
$\mathbb{Z}^{\nu}$ ; la seconde est l'action sur l'ensemble des bords et la
signature.

La question essentielle (page 62) : reconstituer, pour un graphe pondéré à
repères $\underline{\Gamma}$, le groupoïde
$\mathcal{T}^{\pm}_{\mathrm{top}}(\underline{\Gamma})$ des objets de graphe
$\underline{\Gamma}$, fonctoriellement en $\underline{\Gamma}$, avec les sept
opérations sauf $4^\circ$ et $5^\circ$. L'opération $4^\circ$ viendra d'elle-même
avec la description explicite des $\mathcal{T}^{\pm}_{\mathrm{top}}(\underline{\Gamma})$,
faite comme pour les $\mathcal{T}^{\pm}_{g,\nu}$ et $S\mathcal{T}^{\pm}_{g,\nu}$
connus. Les bouchages de trous s'obtiendraient à partir des cas particuliers à
ajouter aux structures « élémentaires », les diminutions
\[
(**)\qquad S_{*}\mathcal{T}^{\pm}_{g,\nu,\nu'} \longrightarrow
S_{*}\mathcal{T}^{\pm}_{g,\nu-\nu'},
\]
où $\mathcal{T}^{\pm}_{g,\nu,\nu'}$ a pour objets les surfaces de type $(g,\nu)$
munies d'un ensemble $I'$ de $\nu'$ bords, qu'on bouche\footnote{Le signe entre
$\nu$ et $\nu'$ n'est pas lu ; boucher $\nu'$ trous donne le type $(g, \nu -
\nu')$, d'où le « $-$ ». La page demande « avec ou sans $S$ ? ».}.

Mais (page 63) ces bouchages sont superficiels : ils ne relient pas les
$\mathcal{T}_{g,\nu}$ aux $\mathcal{T}_{g',\nu}$ de genre $g' < g$. Pour cela il
faut le gommage des cercles de découpage — « ex-cercles de recollement ! » : un
cercle obtenu en recollant deux bords, qu'on efface ensuite\footnote{La page
appelle cette opération $5^\circ$ ; c'est le $4^\circ$ de la récapitulation de
la page 60, et on la nomme ainsi.}.

Autre structure importante à expliciter : les groupoïdes fondamentaux $\Pi_1(X;
R)$ des $(X, \widetilde{\Sigma}, R)$, à points-base dans $R \neq \emptyset$, en
tant que groupoïdes dépendant fonctoriellement de $R$, avec leur comportement
par les opérations — un nouvel objet en groupoïdes au-dessus de
$S_{*}\mathcal{T}^{\pm}_{\mathrm{top}}$. Déjà pour $S_2\mathcal{T}_{0,3}$, un
pantalon avec deux repères, la question est non triviale. C'est sans doute trop
ambitieux : il faudrait un « modèle » de $S_{*}\mathcal{T}^{\pm}_{\mathrm{top}}$
assez gros pour contenir « toutes les cordes — ça fait un peu gros ! ». Mais on
doit pouvoir déterminer sans trop de mal, dans le modèle de
$S_{*}\mathcal{T}D^{\pm}_{\mathrm{top}}$, des modèles des $\Pi_1(X, R)$ et leur
dépendance fonctorielle en $(X, R)$ — et la suite s'arrête sur ces points de
suspension\footnote{La fin de la phrase de la page 63 est chargée de ratures et
d'ajouts, où passent des revêtements finis $X'$ des $X$ ; l'ordre des morceaux
est incertain et on n'en retient que ce qui précède. Un groupe modulaire agissant
sur un groupoïde fondamental à points-base sur le bord, c'est le cadre où le
groupe de Galois de $\overline{\mathbb{Q}}$ agira aussi, par la tour des
$\mathcal{T}_{\overline{\mathbb{Q}}\text{-alg}}$ de la page 42 ;
rapprochement de cette édition.}.

\pagerange{65}{68}
\section*{Découpages en pantalons (pages 65 à 68)}

La page 65, une couverture, porte « Découpage de cartes en “pantalons”
généralisés ».

\emph{Le pantalon pondéré.} Un pantalon subdivisé de façon standard, d'ensemble
de trous $J$ de cardinal $3$, équivaut à un « triangle pondéré » $\pi : J \to
\mathbb{N}$, où $\pi(j)$ est le nombre d'arêtes qui relient les deux trous autres
que $j$\footnote{La page écrit « $\pi(j)$ = nb d'arêtes qui relient $j$ et $k$,
où $i \neq j \neq k$ » ; le dessin (où $\alpha = 3$ compte les points du côté
$1\infty$, opposé au trou $0$) et les relations qui suivent imposent la
lecture donnée.}. Si $\alpha, \beta, \gamma$ sont ces nombres pour les trous
$0, 1, \infty$, le nombre de points sur le cercle de chaque trou est
\[
\alpha' = \beta + \gamma, \qquad \beta' = \gamma + \alpha, \qquad \gamma' =
\alpha + \beta,
\]
d'où
\[
\alpha = \frac{\beta' + \gamma' - \alpha'}{2}, \qquad \beta = \frac{\gamma' +
\alpha' - \beta'}{2}, \qquad \gamma = \frac{\alpha' + \beta' - \gamma'}{2}.
\]
Les $\alpha, \beta, \gamma$ sont entiers si et seulement si $\alpha' + \beta' +
\gamma'$ est pair, et entiers naturels si de plus $\alpha', \beta', \gamma'$
vérifient les inégalités triangulaires\footnote{La page n'énonce que la parité, «
$\in \mathbb{N}$ » lu sur un signe rapide. C'est aujourd'hui la paramétrisation
des arcs dans un pantalon par leurs nombres d'intersection avec les bords, l'un
des ingrédients des coordonnées de Dehn–Thurston ; rapprochement de cette
édition.}. Sur la sphère à trois sommets $0, 1, \infty$ (les trous réduits à des
points), la carte a $\alpha + \beta + \gamma$ arêtes et, si $\alpha, \beta,
\gamma \geq 1$, $(\alpha - 1) + (\beta - 1) + (\gamma - 1) = \alpha + \beta +
\gamma - 3$ faces-bigones et $2$ faces-triangles, soit $\alpha + \beta + \gamma
- 1$ faces, et la formule d'Euler $3 - (\alpha + \beta + \gamma) + (\alpha +
\beta + \gamma - 1) = 2$ est satisfaite\footnote{La page écrit « $= \alpha +
\beta + \gamma - 2$ » pour le nombre de bigones, le $2$ douteux ; c'est $-3$,
qui est aussi ce que donne son propre total de faces.}.

\emph{La carte pantalonnière généralisée.} Une carte orientée découpée en
pantalons de façon standard équivaut à une carte trivalente (avec
éventuellement des sommets d'ordre $1$) dont les couples (sommet, arête
incidente) sont pondérés, de sorte que, pour chaque arête $st$,
\[
\alpha_s + \beta_s = \alpha_t + \beta_t \;(= \gamma),
\]
le nombre de points sur le cercle de recollement vu des deux pantalons ; avec,
de plus, un anti-isomorphisme entre les torseurs sous $\mathbb{Z}/\gamma$ que sont
les sommes circulaires de $[1, \alpha_s]$ et $[1, \beta_s]$, et de $[1, \alpha_t]$
et $[1, \beta_t]$. C'est exactement la donnée de recollement $7^\circ$ de la
page 60, un élément de $R_a \wedge R_{a'}$, pour une décomposition en pantalons
où chaque cercle porte les $\gamma$ points où la carte le traverse\footnote{« de
données d'une anti-isom. » : « anti- » est suppléé ; « $\mathbb{Z}/\gamma$ » est
douteux, et la flèche entre les deux sommes circulaires est rendue par
l'édition.}.

Les pages 66 à 68 dessinent ces découpages : la sphère à trois trous $0, 1,
\infty$ subdivisée, ses points $a_i, b_i$ rangés par trou, des triangles
pondérés ($\alpha = 3$, $\beta = 5$, $\gamma = 4$), une carte trivalente, et,
page 68, une carte sur la sphère dont les sommets sont entourés de courbes
fermées emboîtées, à l'encre rouge — un découpage en pantalons. La page 67 ne
porte que des esquisses de la page 66.

\pagerange{69}{70}
\section*{Une autre mouture des définitions (pages 69 et 70)}

La page 69, sur un décalque du dessin de la page 68, porte un carré
\[
\begin{array}{ccc}
\mathcal{T}^{\pm}_{\mathrm{top}} & \xrightarrow{\ \mathrm{imm}\ } &
S_{*}\mathcal{T}^{\pm}_{\mathrm{top}} \\
\uparrow & & \uparrow \\
\mathcal{T}^{+}_{\mathrm{top}} & \xrightarrow{\ \mathrm{imm}\ } &
S_{*}\mathcal{T}^{+}_{\mathrm{top}}
\end{array}
\]
au-dessus de deux $[\pm 1]$ : les flèches horizontales sont des immersions
ouvertes (inclusions de sous-groupoïdes pleins réunions de composantes), les
verticales sont fidèles et essentiellement surjectives, sans être pleines ; d'où
$\pi_0(\mathcal{T}^{+}_{\mathrm{top}}) \xrightarrow{\sim}
\pi_0(\mathcal{T}^{\pm}_{\mathrm{top}})$, la ligne s'arrêtant là\footnote{La page
écrit $\mathcal{T}S_{\mathrm{top}}$ et $S\mathcal{T}S^{\pm}_{\mathrm{top}}$, le
premier $S$ lu sous réserve ; d'après la page 70, ce sont les groupoïdes
spéciaux.}.

La page 70, intitulée elle aussi « Groupoïdes de Teichmüller », recommence les
définitions des pages 43 à 45 dans un autre ordre et avec d'autres
lettres\footnote{La page note $S$ l'ensemble des points marqués, qu'on écrit
$R$ ; son préfixe $S$, sans étoile, est le $S_{*}$ des pages 43 sqq.}.
\begin{enumerate}
\item[1)] $S_{*}\mathcal{T}^{+}_{\mathrm{top}}$, le groupoïde de Teichmüller
\emph{spécial} orienté topologique : objets les paires $(X, R)$, $X$ surface
compacte orientée à bord, $R \subset \partial X$ finie — les repères sont ici sur
le bord seulement ; flèches les classes d'isotopie d'isomorphismes. Pour $R =
\emptyset$, on a $\mathcal{T}^{+}_{\mathrm{top}}$, sous-groupoïde ouvert. Variante
$S_{*}\mathcal{T}^{\pm}_{\mathrm{top}}$, avec $\varepsilon :
S_{*}\mathcal{T}^{\pm}_{\mathrm{top}} \to [\pm 1]$ (l'orientation), de fibre
$S_{*}\mathcal{T}^{+}_{\mathrm{top}}$ — sur les surfaces connexes, comme à la
page 43.
\item[2°)] Les triples $(X, \Sigma, R)$, $\Sigma$ sous-variété compacte de
dimension $1$ de $X \setminus \partial X$, $R \subset \Sigma \cup \partial X$
finie : $\mathcal{T}D^{+}_{\mathrm{top}}$, « $D$ initiale de \emph{divisage} »,
groupoïde de Teichmüller orienté divisé ; et $S_{*}\mathcal{T}^{+}_{\mathrm{top}}
\subset S_{*}\mathcal{T}D^{+}_{\mathrm{top}}$ est l'immersion ouverte des objets
où $\Sigma = \emptyset$.
\item[3°)] De même le biorienté divisé.
\end{enumerate}
La différence avec les pages 43 à 45 n'est pas que de lettres : ici les repères
de $S_{*}\mathcal{T}^{+}_{\mathrm{top}}$ sont sur le bord, et le $\Sigma$ de 2°)
est disjoint du bord, alors que la page 45 met $\partial X$ dans
$\widetilde{\Sigma}$ ; les deux présentations décrivent les mêmes objets.

\pagerange{71}{71}
\section*{La page 71}

Une dernière planche, légendée « découpage pantalonnesque d'une carte
(triangulée) » : une triangulation de neuf sommets, neuf courbes fermées
emboîtées qui la découpent en pantalons, et le graphe trivalent pondéré
correspondant, dont les poids sont inscrits aux sommets et aux arêtes — la
donnée de la page 66 sur un exemple.

\end{document}
